\providecommand{\Ad}{\operatorname{Ad}}
\providecommand{\Hol}{\operatorname{Hol}}
\providecommand{\GL}{\operatorname{GL}}
\providecommand{\SO}{\operatorname{SO}}
\providecommand{\SU}{\operatorname{SU}}
\providecommand{\Ut}{\operatorname{U}}
\providecommand{\Ot}{\operatorname{O}}
\providecommand{\Mani}{\mathcal{M}}
\providecommand{\Lie}[1]{\mathfrak{#1}}
\providecommand{\Cl}{\operatorname{Cl}}
\providecommand{\Stab}{\operatorname{Stab}}
\providecommand{\Img}{\operatorname{Im}}
\providecommand{\rank}{\operatorname{rank}}
\providecommand{\dR}{\mathrm{dR}}
\providecommand{\softmax}{\operatorname{softmax}}

\chapter{Gauges: redes gauge-equivariantes}
\label{ch:gauges}

El cap\'itulo anterior dot\'o a las variedades de geometr\'ia intr\'inseca --- una
m\'etrica, geod\'esicas, una conexi\'on --- y la us\'o para transportar
caracter\'isticas y construir convoluciones en dominios curvos
(\Cref{ch:geodesics}). Una obstrucci\'on sutil qued\'o all\'i impl\'icita: para
escribir un vector, un filtro o cualquier cantidad direccional en una variedad,
primero hay que elegir una base del espacio tangente en cada punto, un
\emph{marco de referencia local} o \emph{gauge}. En un dominio curvo esta
elecci\'on es inevitable, no puede hacerse de forma consistente y global, y es
completamente arbitraria. Dos investigadores que analizan la misma superficie
pueden escoger marcos distintos en cada punto y obtener mapas de
caracter\'isticas num\'ericamente diferentes que describen exactamente la misma
geometr\'ia.

La teor\'ia gauge es el lenguaje matem\'atico de esta libertad. Tomada de la
f\'isica te\'orica, donde las simetr\'ias gauge son las transformaciones internas
que dejan invariantes las ecuaciones de un sistema, proporciona las herramientas
para construir redes neuronales cuyas predicciones son independientes de la
elecci\'on arbitraria del marco local. Esta es la quinta y m\'as general de las
cinco G: donde las cuadr\'iculas (grids) explotan una simetr\'ia global de
traslaci\'on y los grupos (groups) un grupo de simetr\'ia global, los gauges
tratan simetr\'ias que act\'uan \emph{localmente}, variando de un punto a otro. La
recompensa es un principio unificador. Las redes equivariantes a grupos, las
redes esf\'ericas y las convoluciones en variedades resultan todas ser casos
particulares de la convoluci\'on gauge-equivariante sobre un fibrado principal
apropiado.

Este cap\'itulo desarrolla el marco en varios movimientos. Primero, el andamiaje
geom\'etrico: los fibrados principales, que registran de golpe todos los marcos
locales posibles, y las conexiones que nos permiten comparar marcos en puntos
distintos (\Cref{sec:gauge:bundles,sec:gauge:transport}). Segundo, la simetr\'ia
en s\'i: las transformaciones gauge, los campos de caracter\'isticas
gauge-equivariantes y la curvatura como obstrucci\'on a una comparaci\'on
independiente del camino
(\Cref{sec:gauge:symmetry,sec:gauge:curvature}). Tercero, la construcci\'on de la
convoluci\'on gauge-equivariante y la red convolucional icosa\'edrica de
\citet{cohen2019gaugeequivariantconvolutional} como realizaci\'on concreta y
eficiente (\Cref{sec:gauge:conv}). Por \'ultimo, repasamos las generalizaciones
algebraicas y topol\'ogicas que abre la perspectiva gauge --- \'algebras de
Clifford y geom\'etricas, complejos simpliciales y celulares, laplacianos de
Hodge, homolog\'ia persistente y las correspondientes arquitecturas de
atenci\'on (\Cref{sec:gauge:clifford,sec:gauge:topology}) --- junto con las
codificaciones posicionales, el rompimiento controlado de simetr\'ia y la
s\'intesis de paso de mensajes que enlaza todo el libro
(\Cref{sec:gauge:positional,sec:gauge:synthesis}). A lo largo del cap\'itulo nos
apoyamos en las conexiones y el transporte paralelo introducidos en
\Cref{ch:geodesics}, reformul\'andolos en el lenguaje m\'as general de fibrados
que exige la simetr\'ia gauge.

\section{De la simetr\'ia global a la local}
\label{sec:gauge:local}

Las redes convolucionales cl\'asicas suponen que la misma transformaci\'on ---
la traslaci\'on --- es v\'alida en todas partes del dominio. Esto funciona para
im\'agenes planas pero se rompe sistem\'aticamente en cuanto los datos llevan
estructura geom\'etrica intr\'inseca.

\begin{example}[Rotar una esfera]
\label{ex:gauge:sphere-rotation}
Consideremos una imagen de la Tierra pintada sobre una esfera y una rotaci\'on
global de $90^\circ$ alrededor del eje polar. En el ecuador la rotaci\'on
preserva fielmente las formas locales; cerca de los polos la misma rotaci\'on las
distorsiona severamente, y un continente cercano al polo norte se ve
completamente distinto despu\'es. La transformaci\'on apropiada depende de la
posici\'on: una \'unica simetr\'ia global es el objeto equivocado en un dominio
curvo.
\end{example}

En un dominio curvo, aplicar una misma transformaci\'on en todas partes puede
romper la estructura geom\'etrica natural, introducir distorsiones artificiales
ausentes en los datos y violar los principios f\'isicos que rigen el fen\'omeno
modelado. El remedio es permitir que la transformaci\'on var\'ie de un punto a
otro.

La analog\'ia f\'isica es instructiva. En electromagnetismo los campos
observables $\mathbf{E}$ y $\mathbf{B}$ no dependen de c\'omo describamos el
potencial electromagn\'etico $A_\mu(x)$: podemos transformar
$A_\mu(x) \mapsto A_\mu(x) + \partial_\mu \Lambda(x)$ con $\Lambda$ una funci\'on
arbitraria que var\'ia de un punto a otro, y los campos
\[
  \mathbf{E} = -\nabla \phi - \frac{\partial \mathbf{A}}{\partial t},
  \qquad
  \mathbf{B} = \nabla \times \mathbf{A}
\]
permanecen invariantes. Esta es la simetr\'ia \emph{local} (gauge) prototípica.
As\'i como las leyes de la f\'isica no deben depender de nuestra elecci\'on de
descripci\'on local, las representaciones neuronales no deber\'ian depender de
elecciones arbitrarias del marco de referencia local.

Esto sugiere una progresi\'on natural en generalidad:
\begin{enumerate}
  \item \emph{Redes cl\'asicas}: una \'unica transformaci\'on (traslaci\'on)
    v\'alida globalmente;
  \item \emph{Redes equivariantes a grupos}: un grupo finito fijo de
    transformaciones v\'alido globalmente (\Cref{ch:groups});
  \item \emph{Convoluciones en variedades}: transformaciones dictadas por la
    geometr\'ia espec\'ifica de la variedad (\Cref{ch:geodesics});
  \item \emph{Redes gauge-equivariantes}: libertad total para elegir
    transformaciones locales, sujeta a un requisito de consistencia.
\end{enumerate}

Se podr\'ia objetar: \textquestiondown por qu\'e no aumentar simplemente los datos con muchas
transformaciones locales? Porque el aumento de datos espera lograr la invarianza
\emph{estad\'isticamente}, a trav\'es de muchas muestras y muchas pasadas hacia
adelante, mientras que la equivarianza gauge garantiza la invarianza
\emph{exacta} por construcci\'on, en una sola pasada hacia adelante. M\'as
all\'a de la eficiencia, la equivarianza exacta impide que la red se aferre a
orientaciones accidentales que no reflejan la verdadera estructura del problema.

\section{Marcos de referencia locales y fibrados principales}
\label{sec:gauge:bundles}

El punto de partida es la estructura de un espacio homog\'eneo $G/H$ estudiada en
\Cref{ch:foundations}. Espacios como la esfera $S^{n-1} \cong \SO(n)/\SO(n-1)$
surgen como cocientes $G/H$, e identificar el estabilizador de un punto $p$ con
el subgrupo $H$ requiere elegir un representante $g$ de la clase lateral $gH$.
Cualquier otro representante $g' = gh$ con $h \in H$ describe el mismo punto pero
da una identificaci\'on distinta. Esta arbitrariedad es la \emph{libertad gauge},
y en general no hay forma can\'onica de fijar la elecci\'on de manera consistente
sobre toda la variedad.

La soluci\'on matem\'atica es no elegir. En lugar de escoger un representante
arbitrario en cada punto, un \emph{fibrado principal} registra \emph{todas} las
elecciones posibles simult\'aneamente, adjuntando a cada punto de la variedad
base el grupo entero de marcos de referencia posibles.

\begin{definition}[Fibrado principal]
\label{def:gauge:principal-bundle}
Un \emph{fibrado principal} con grupo de estructura $G$ es una c\'uadrupla
$(P, \Mani, \pi, G)$ donde $P$ (el \emph{espacio total}) y $\Mani$ (la
\emph{base}) son variedades suaves, $\pi\colon P \to \Mani$ es una submersi\'on
sobreyectiva (la \emph{proyecci\'on}), y $G$ es un grupo de Lie que act\'ua
\emph{libremente} sobre $P$ por la derecha --- es decir, $u \cdot g = u$ para
alg\'un $u$ obliga a $g = e$, de modo que $\Stab(u) = \{e\}$ para todo $u$ ---
sujeto a:
\begin{enumerate}
  \item[(a)] para cada $p \in \Mani$ la fibra $\pi^{-1}(p)$ es difeomorfa a
    $G$, la \emph{fibra t\'ipica}, de modo que $\dim P = \dim \Mani + \dim G$; el
    difeomorfismo $\pi^{-1}(p) \cong G$ no es can\'onico, y fijarlo equivale a
    elegir un gauge en $p$;
  \item[(b)] la acci\'on preserva las fibras, $\pi(u \cdot g) = \pi(u)$;
    combinada con la libertad, $G$ act\'ua libre y transitivamente sobre cada
    fibra, convirti\'endola en un \emph{torsor} de $G$ --- una copia de $G$ sin
    identidad distinguida;
  \item[(c)] existen \emph{trivializaciones locales}: todo $p$ tiene un
    entorno $U$ y un difeomorfismo $\phi_U\colon \pi^{-1}(U) \to U \times
    G$ con $\phi_U(u \cdot g) = (\pi(u), \phi_2(u)\, g)$. En los solapamientos,
    \emph{funciones de transici\'on} $g_{UV}\colon U \cap V \to G$, definidas por
    $\phi_V \circ \phi_U^{-1}(p, g) = (p, g_{UV}(p)\, g)$, relacionan dos
    trivializaciones y codifican completamente la topolog\'ia global del fibrado.
\end{enumerate}
Para un desarrollo sistem\'atico v\'ease~\citep{Kobayashi1996}.
\end{definition}

\begin{remark}[Las fibras como subvariedades]
\label{rem:gauge:fibres-submanifold}
Como $\pi$ es una submersi\'on, el teorema del conjunto de nivel regular
garantiza que cada fibra $\pi^{-1}(p)$ es una subvariedad suave de $P$ de
dimensi\'on $\dim P - \dim \Mani = \dim G$. Esto justifica la condici\'on (a):
siendo subvariedades de la dimensi\'on correcta sobre las que $G$ act\'ua libre y
transitivamente, el difeomorfismo $\pi^{-1}(p) \cong G$ queda fijado al elegir
cualquier punto \'unico $u_0 \in \pi^{-1}(p)$, mediante $g \mapsto u_0 \cdot g$.
\end{remark}

La libertad codificada en una fibra es precisamente la libertad de elegir un
elemento dentro de una clase lateral $gH$. El v\'inculo entre las descripciones
algebraica y geom\'etrica es directo.

\begin{proposition}[Los espacios homog\'eneos son fibrados principales]
\label{prop:gauge:homog-bundle}
Sea $G$ un grupo de Lie y $H \le G$ un subgrupo cerrado. Entonces
$(G, G/H, \pi, H)$, con $\pi$ la proyecci\'on cociente y $H$ actuando por
multiplicaci\'on a la derecha, es un fibrado principal.
\end{proposition}

\begin{proof}
Comprobamos las condiciones de \Cref{def:gauge:principal-bundle}.
\emph{(a)} Para una clase lateral $gH \in G/H$ la fibra es $\pi^{-1}(gH) = gH$, y
la traslaci\'on izquierda $L_g\colon H \to gH$, $h \mapsto gh$, es un
difeomorfismo; por tanto cada fibra es difeomorfa a $H$.
\emph{(b)} La acci\'on derecha de $H$ preserva las fibras: para $g \in G$,
$h \in H$, $\pi(gh) = (gh)H = gH = \pi(g)$ ya que $hH = H$.
\emph{(c)} Las trivializaciones locales son equivalentes a secciones suaves
locales $s\colon U \to G$ de $\pi$, que existen para cualquier subgrupo cerrado
$H$ de un grupo de Lie~\citep[Ch.~I, Prop.~5.5]{Kobayashi1996}.
\end{proof}

As\'i, los espacios homog\'eneos obtenidos mediante el teorema
\'orbita--estabilizador de \Cref{ch:foundations} son autom\'aticamente bases de
fibrados principales: $S^{n-1} \cong \SO(n)/\SO(n-1)$, $\mathbb{H}^n \cong
\SO^+(n,1)/\SO(n)$, y as\'i sucesivamente. En cada caso el estabilizador $H$ se
convierte en la \emph{fibra} (las simetr\'ias locales que fijan un punto del
dominio), mientras que la \'orbita $G/H$ es la \emph{base} sobre la que viven las
se\~nales. Las libertades gauge formalizadas en este cap\'itulo son precisamente
las libertades de elegir un elemento dentro de cada clase lateral $gH$.
\Cref{fig:gauge:bundle} ilustra la estructura.

\begin{figure}[ht]
\centering
\begin{tikzpicture}[>=Stealth, font=\small]
% Total space P (top)
\fill[gdlgreen!12, rounded corners=6pt] (-2.0, 1.0) rectangle (2.4, 4.2);
\draw[gdlgreen!50!black, dashed, rounded corners=6pt] (-2.0, 1.0) rectangle (2.4, 4.2);
\node[gdlgreen!50!black, anchor=south] at (0.2, 4.2) {$\pi^{-1}(U)$};
\draw[thick, gdlblue!40!black, rounded corners=8pt, dashed] (-4.4, 0.8) rectangle (3.0, 4.7);
\node[gdlblue!40!black, anchor=south] at (-0.7, 4.7) {$P$};
% fibre outside U
\draw[thick, gdlorange!60, fill=gdlorange!15] (-3.4, 2.6) ellipse (0.28cm and 0.75cm);
\foreach \yy/\op in {3.1/65, 2.6/50, 2.1/35} \fill[gdlred!\op] (-3.4, \yy) circle (1.5pt);
\draw[gdlgray!40, dashed] (-3.4, 1.85) -- (-3.4, -0.2);
\node[gdlorange!70!black, left] at (-3.7, 2.6) {$G$};
% fibres over U
\foreach \xx in {-1.2, 0.2, 1.6} {
  \draw[thick, gdlorange!70, fill=gdlorange!18] (\xx, 2.6) ellipse (0.28cm and 0.75cm);
  \foreach \yy/\op in {3.1/65, 2.6/50, 2.1/35} \fill[gdlred!\op] (\xx, \yy) circle (1.5pt);
  \draw[gdlgray!40, dashed] (\xx, 1.85) -- (\xx, -0.2);
}
\node[gdlorange!70!black, right] at (1.95, 3.2) {$\pi^{-1}(p)\cong G$};
% projection
\draw[->, very thick, gdlpurple!70] (0, 0.7) -- node[right] {$\pi$} (0, 0.2);
% base M
\fill[gdlblue!15] (-4.4, -0.6) .. controls (-3, 0.0) and (-1, -0.8) .. (0.5, -0.5)
   .. controls (2, -0.2) and (2.4, -0.7) .. (3.0, -0.55)
   -- (3.0, -1.4) .. controls (2.4, -1.5) and (2, -1.2) .. (0.5, -1.5)
   .. controls (-1, -1.8) and (-3, -1.0) .. (-4.4, -1.6) -- cycle;
\fill[gdlgreen!22, rounded corners=4pt] (-2.0, -0.55) rectangle (2.4, -1.35);
\node[gdlgreen!50!black] at (0.2, -0.95) {$U$};
\node[gdlblue!70!black] at (-3.7, -0.3) {$\Mani$};
\foreach \xx in {-1.2, 0.2, 1.6} \fill[gdlgreen!50!black] (\xx, -0.75) circle (1.5pt);
\node[gdlgreen!50!black, below, font=\tiny] at (0.2, -0.85) {$p$};
\fill[gdlblue!60] (-3.4, -0.75) circle (1.5pt);
% product side
\begin{scope}[shift={(7.0,0.0)}]
  \fill[gdlgreen!15] (-1.0, 0.0) rectangle (2.0, 3.2);
  \draw[thick, gdlgreen!50!black] (-1.0, 0.0) rectangle (2.0, 3.2);
  \draw[->, thick] (-1.0, 0.0) -- (2.3, 0.0); \node[below] at (0.6,-0.05) {$U$};
  \draw[->, thick] (-1.0, 0.0) -- (-1.0, 3.5); \node[left] at (-1.0,1.6) {$G$};
  \node[gdlgreen!40!black] at (0.6, 3.5) {$U \times G$};
  \foreach \xpos in {-0.3, 0.6, 1.5} {
    \draw[gdlorange!60, thick] (\xpos, 0.1) -- (\xpos, 3.1);
    \foreach \yy/\op in {2.4/65,1.6/50,0.8/35} \fill[gdlred!\op] (\xpos, \yy) circle (1.5pt);
  }
\end{scope}
\draw[<->, thick, gdlpurple!80] (2.6, 2.6) -- node[above] {$\phi_U$} node[below, font=\scriptsize, gdlpurple!60] {$\cong$} (5.9, 1.6);
\end{tikzpicture}
\caption[Fibrado principal]{Un fibrado principal y su trivializaci\'on local. El
espacio total $P$ se proyecta sobre la variedad base $\Mani$ mediante $\pi$; cada
fibra $\pi^{-1}(p)$ es difeomorfa al grupo de estructura $G$ y re\'une todos los
marcos locales posibles en $p$. Sobre un entorno $U$ el difeomorfismo
$\phi_U$ presenta el fibrado como el producto $U \times G$, pero su topolog\'ia
global puede estar retorcida.}
\label{fig:gauge:bundle}
\end{figure}

\begin{example}[Fibrado de marcos ortonormales]
\label{ex:gauge:frame-bundle}
Sea $(\Mani, g)$ una variedad riemanniana de dimensi\'on $n$. Su
\emph{fibrado de marcos ortonormales} $\mathrm{FM}(\Mani)$ es el fibrado
principal con grupo de estructura $\Ot(n)$ cuyo espacio total consiste en todos
los marcos ortonormales $(p, e_1, \ldots, e_n)$ con $\{e_i\}$ una base ortonormal
de $T_p\Mani$; la proyecci\'on $(p, e_1, \ldots, e_n) \mapsto p$ olvida el marco,
y $A \in \Ot(n)$ act\'ua por $(p, e_1, \ldots, e_n) \cdot A = (p, \sum_j A_{j1}
e_j, \ldots, \sum_j A_{jn} e_j)$. Cada fibra re\'une toda elecci\'on posible de
marco ortonormal en $p$ --- precisamente la libertad gauge local. Para una
superficie ($n = 2$) la libertad rotacional relevante es $\SO(2)$.
\end{example}

\begin{example}[Fibrado trivial]
\label{ex:gauge:trivial-bundle}
El fibrado m\'as simple es el producto $P = \Mani \times G$ con $\pi(p, g) = p$ y
acci\'on $(p, h) \cdot g = (p, hg)$. Admite una secci\'on global $s(p) = (p, e)$
--- una elecci\'on de gauge globalmente consistente, con $e$ la identidad. No
todo fibrado admite una secci\'on global; si lo hace o no es una cuesti\'on
topol\'ogica, retomada en \Cref{sec:gauge:curvature}.
\end{example}

En el lenguaje del aprendizaje profundo geom\'etrico, la base $\Mani$ es donde
viven los datos (por ejemplo la esfera $S^2$ para se\~nales esf\'ericas),
mientras que el grupo de estructura $G$ codifica las simetr\'ias locales
relevantes ($\SO(2)$ para rotaciones de un marco tangente). La fibra sobre cada
punto contiene todas las orientaciones o marcos de referencia locales posibles
all\'i.

\section{Conexiones y transporte paralelo}
\label{sec:gauge:transport}

Una conexi\'on en un fibrado principal prescribe c\'omo transportar informaci\'on
entre fibras distintas de un modo compatible con la estructura de grupo. Divide
el espacio tangente del espacio total en una parte \emph{vertical} a lo largo de
la fibra y una parte \emph{horizontal} transversal a ella. Esta comparaci\'on es
lo que hace que operaciones como la convoluci\'on sean definibles en dominios no
eucl\'ideos.

\begin{definition}[Espacio vertical y campo vectorial fundamental]
\label{def:gauge:vertical}
Para $u \in P$ el \emph{espacio vertical} es
$V_u P = \ker(\mathrm{d}\pi_u) = \{X \in T_u P : \mathrm{d}\pi_u(X) = 0\}$, las
direcciones tangentes a la fibra. Para $A \in \Lie{g}$, el \'algebra de Lie de
$G$, el \emph{campo vectorial fundamental} generado por $A$ en $u$ es
\[
  A^*_u = \left.\frac{\mathrm{d}}{\mathrm{d}t}\right|_{t=0} u \cdot \exp(tA),
\]
la velocidad de la curva de \'orbita $t \mapsto u \cdot \exp(tA)$. Como esta
curva permanece en la fibra $\pi^{-1}(\pi(u))$, tenemos $A^*_u \in V_u P$.
\end{definition}

\begin{proposition}[Isomorfismo can\'onico $V_u P \cong \Lie{g}$]
\label{prop:gauge:vertical-iso}
Para cada $u \in P$ la aplicaci\'on $\sigma_u\colon \Lie{g} \to V_u P$,
$\sigma_u(A) = A^*_u$, es un isomorfismo lineal.
\end{proposition}

\begin{proof}
Consideremos la aplicaci\'on de \'orbita $\Phi_u\colon G \to P$, $\Phi_u(g) = u
\cdot g$. Su diferencial en la identidad es $\sigma_u = (\mathrm{d}\Phi_u)_e\colon
T_e G = \Lie{g} \to T_u P$, con imagen dentro de $V_u P$ por
\Cref{def:gauge:vertical}; siendo el diferencial de una aplicaci\'on suave es
lineal. Para la inyectividad, la libertad de la acci\'on hace que $\Phi_u$ sea
inyectiva y una inmersi\'on: si $(\mathrm{d}\Phi_u)_e(A) = 0$ la curva
$t \mapsto u \cdot \exp(tA)$ tiene velocidad inicial nula, lo que por la libertad
fuerza $A = 0$. Para la sobreyectividad, como $\pi$ es una submersi\'on,
$\dim V_u P = \dim P - \dim \Mani = \dim G = \dim \Lie{g}$, y una aplicaci\'on
lineal inyectiva entre espacios de igual dimensi\'on finita es un isomorfismo.
\end{proof}

\begin{definition}[Conexi\'on]
\label{def:gauge:connection}
Una \emph{conexi\'on} en $(P, \Mani, \pi, G)$ es una $1$-forma con valores en el
\'algebra de Lie $\omega \in \Omega^1(P, \Lie{g})$ tal que
\begin{enumerate}
  \item \emph{normalizaci\'on}: $\omega(A^*_u) = A$ para todo $A \in \Lie{g}$ y
    $u \in P$;
  \item \emph{equivarianza}: $R_g^* \omega = \Ad_{g^{-1}} \omega$ para todo
    $g \in G$, donde $R_g$ es la acci\'on derecha y $\Ad$ la representaci\'on
    adjunta.
\end{enumerate}
Equivalentemente, una conexi\'on define una \emph{distribuci\'on horizontal}
$H_u P = \ker(\omega_u) \subset T_u P$ complementaria al espacio vertical, de modo
que $T_u P = H_u P \oplus V_u P$ y $H_{u \cdot g} P = (R_g)_* H_u P$.
\end{definition}

La distribuci\'on horizontal especifica qu\'e direcciones del espacio total son
``puramente espaciales'' (paralelas a la base) frente a ``puramente gauge'' (a lo
largo de la fibra). Una curva $\tilde\gamma$ en $P$ es \emph{horizontal} si
$\dot{\tilde\gamma}(t) \in H_{\tilde\gamma(t)} P$ para todo $t$; tales curvas
representan el transporte del marco ``sin rotaci\'on''. Esta es exactamente la
forma en lenguaje de fibrados de la conexi\'on af\'in de \Cref{ch:geodesics}.

\begin{definition}[Transporte paralelo]
\label{def:gauge:parallel-transport}
Dada una conexi\'on y un camino suave $\gamma\colon [0,1] \to \Mani$, el
\emph{transporte paralelo} es la aplicaci\'on
$\Gamma_\gamma\colon \pi^{-1}(\gamma(0)) \to \pi^{-1}(\gamma(1))$ definida como
sigue: para $u_0 \in \pi^{-1}(\gamma(0))$ existe un \'unico \emph{levantamiento
horizontal} $\tilde\gamma\colon [0,1] \to P$ con (i) $\pi \circ \tilde\gamma =
\gamma$, (ii) $\tilde\gamma(0) = u_0$, y (iii) $\dot{\tilde\gamma}(t) \in
H_{\tilde\gamma(t)} P$ para todo $t$; entonces $\Gamma_\gamma(u_0) =
\tilde\gamma(1)$.
\end{definition}

\begin{theorem}[Propiedades del transporte paralelo]
\label{thm:gauge:transport-props}
El transporte paralelo satisface:
\begin{enumerate}
  \item \emph{equivarianza}: $\Gamma_\gamma(u \cdot g) = \Gamma_\gamma(u) \cdot
    g$ para todo $g \in G$;
  \item \emph{composici\'on}: $\Gamma_{\gamma_2 * \gamma_1} = \Gamma_{\gamma_2}
    \circ \Gamma_{\gamma_1}$ para caminos concatenados;
  \item \emph{inversi\'on}: $\Gamma_{\gamma^{-1}} = (\Gamma_\gamma)^{-1}$, donde
    $\gamma^{-1}(t) = \gamma(1-t)$.
\end{enumerate}
\end{theorem}

\begin{proof}
\emph{Equivarianza.} Sea $\tilde\gamma$ el levantamiento horizontal desde $u$ y
sea $\tilde\gamma_g(t) = \tilde\gamma(t) \cdot g$. Entonces $\pi(\tilde\gamma_g(t))
= \pi(\tilde\gamma(t)) = \gamma(t)$; $\tilde\gamma_g(0) = u \cdot g$; y
$\dot{\tilde\gamma}_g(t) = (R_g)_* \dot{\tilde\gamma}(t) \in (R_g)_*
H_{\tilde\gamma(t)} P = H_{\tilde\gamma(t)\cdot g} P$ por la $G$-invarianza de la
distribuci\'on horizontal. Por la unicidad del levantamiento, $\tilde\gamma_g$ es
el levantamiento desde $u \cdot g$, de modo que $\Gamma_\gamma(u \cdot g) =
\tilde\gamma(1)\cdot g = \Gamma_\gamma(u)\cdot g$.

\emph{Composici\'on.} Sea $\tilde\gamma_1$ el levantamiento horizontal de
$\gamma_1$ desde $u_0$, que termina en $u_1 = \tilde\gamma_1(1)$, y
$\tilde\gamma_2$ el levantamiento de $\gamma_2$ desde $u_1$. La concatenaci\'on
$\tilde\gamma$ se proyecta a $\gamma_2 * \gamma_1$, comienza en $u_0$ y es
horizontal en cada tramo; por unicidad es el levantamiento de $\gamma_2 *
\gamma_1$ desde $u_0$, con punto final $\tilde\gamma_2(1) =
\Gamma_{\gamma_2}(\Gamma_{\gamma_1}(u_0))$.

\emph{Inversi\'on.} Con $\tilde\gamma^-(t) = \tilde\gamma(1-t)$ tenemos
$\pi(\tilde\gamma^-(t)) = \gamma^{-1}(t)$, $\tilde\gamma^-(0) = u_1$, y
$\dot{\tilde\gamma}^-(t) = -\dot{\tilde\gamma}(1-t) \in H_{\tilde\gamma(1-t)} P$,
ya que los subespacios horizontales son espacios vectoriales. Por unicidad este
es el levantamiento de $\gamma^{-1}$ desde $u_1$, dando
$\Gamma_{\gamma^{-1}}(u_1) = u_0 =
\Gamma_\gamma^{-1}(u_1)$~\citep[Ch.~II, Prop.~3.2]{Kobayashi1996}.
\end{proof}

Trabajando en una secci\'on local, el levantamiento horizontal satisface
$\omega(\dot{\tilde\gamma}) = 0$, lo que expresado a trav\'es de una curva con
valores en el grupo $g(t)$ se convierte en la ecuaci\'on diferencial ordinaria
lineal
\[
  \frac{\mathrm{d}g}{\mathrm{d}t} + \omega(\dot\gamma(t))\, g(t) = 0,
\]
con la exponencial ordenada por camino
$g(1) = \mathcal{P}\exp\!\big(-\!\int_0^1 \omega(\dot\gamma)\,\mathrm{d}t\big)\,g(0)$
como soluci\'on formal. Este es el an\'alogo continuo de acumular una matriz de
transporte arista por arista, la receta que implementar\'a una red discreta.

\Cref{fig:gauge:horizontal-vertical} representa la divisi\'on de $T_u P$ y el
papel de la proyecci\'on $\mathrm{d}\pi_u$.

\begin{figure}[ht]
\centering
\begin{tikzpicture}[>=Stealth, font=\small]
% base
\begin{scope}[shift={(0,-2)}]
  \draw[thick, fill=gdlblue!18]
    (-3,0) .. controls (-1,0.5) and (1,-0.3) .. (3,0.2) --
    (3,-0.5) .. controls (1,-1) and (-1,0) .. (-3,-0.5) -- cycle;
  \node[below] at (0,-0.8) {\textbf{base $\Mani$}};
  \fill[gdlblue!70] (0,0.1) circle (3pt);
  \node[below, font=\small] at (0,-0.1) {$x$};
\end{scope}
% fibre
\draw[thick, gdlorange!70, fill=gdlorange!18] (0,0) ellipse (0.4cm and 2cm);
\node[gdlorange!70, right, font=\small] at (0.5,1.5) {fibra $\pi^{-1}(x)\cong G$};
\fill[gdlred!70] (0,0.5) circle (4pt);
\node[right, font=\small, gdlred!70] at (0.2,0.5) {$u$};
% tangent decomposition
\begin{scope}[shift={(0,0.5)}]
  \draw[thick, dashed, gdlgray!40] (-2,-1) -- (2,1) -- (2.5,2) -- (-1.5,0) -- cycle;
  \node[gdlgray!60, font=\small] at (2.5,0) {$T_u P$};
  \draw[->, very thick, gdlpurple!70] (0,0) -- (0,1.2) node[above, font=\small] {$V_u P$};
  \node[gdlpurple!70, font=\tiny, left] at (-0.2,0.6) {vertical};
  \draw[->, very thick, gdlgreen!60!black] (0,0) -- (1.5,0.5) node[right, font=\small] {$H_u P$};
  \node[gdlgreen!60!black, font=\tiny, below] at (0.8,0.1) {horizontal};
  \draw[thick, dashed, gdlblue!50] (0,0) -- (1.5,1.7);
  \node[gdlblue!70, font=\tiny] at (1.0,1.3) {$v = v^V + v^H$};
\end{scope}
\draw[->, thick, gdlgray!70] (1.5,0.8) -- node[right, font=\small] {$\mathrm{d}\pi_u$} (1.5,-1.5);
\node[gdlgray!70, font=\tiny] at (2.7,-0.5) {$\ker(\mathrm{d}\pi_u)=V_u P$};
\node[font=\small, text width=5.2cm] at (6.6,0.2) {
  \textbf{Descomposici\'on:}\\ $T_u P = V_u P \oplus H_u P$\\[0.4em]
  \textbf{Vertical:} $V_u P = \ker(\mathrm{d}\pi_u)$, tangente a la fibra
  (direcciones ``gauge'')\\[0.4em]
  \textbf{Horizontal:} $H_u P$ fijada por la conexi\'on (direcciones
  ``espaciales'')\\[0.4em]
  \textbf{$G$-invarianza:} $H_{u\cdot g} = (R_g)_* H_u$
};
\end{tikzpicture}
\caption[Descomposici\'on horizontal y vertical]{La conexi\'on divide cada
espacio tangente $T_u P$ en una parte vertical $V_u P = \ker(\mathrm{d}\pi_u)$,
tangente a la fibra, y un complemento horizontal $H_u P$. La proyecci\'on
$\mathrm{d}\pi_u$ anula las direcciones verticales; las direcciones horizontales
se proyectan isom\'orficamente sobre la base.}
\label{fig:gauge:horizontal-vertical}
\end{figure}

\section{Simetr\'ia gauge, campos y equivarianza}
\label{sec:gauge:symmetry}

Ahora podemos precisar la simetr\'ia central. Una elecci\'on de gauge es una
elecci\'on de marco en cada punto; cambiarla es una transformaci\'on gauge.

\begin{definition}[Transformaci\'on gauge]
\label{def:gauge:transformation}
Una \emph{transformaci\'on gauge} es una aplicaci\'on suave $g\colon \Mani \to G$
que asigna a cada punto $p$ un elemento $g(p)$ del grupo de estructura. El
conjunto de todas estas aplicaciones forma el \emph{grupo gauge}
\[
  \mathcal{G} = C^\infty(\Mani, G)
\]
bajo la multiplicaci\'on punto a punto, $(g_1 \cdot g_2)(p) = g_1(p)\, g_2(p)$.
\end{definition}

Una transformaci\'on gauge rota el marco local de forma independiente en cada
punto. \Cref{fig:gauge:freedom} muestra dos gauges sobre la misma superficie: los
mismos datos geom\'etricos, descritos por componentes num\'ericas distintas.

\begin{figure}[ht]
\centering
\begin{tikzpicture}[>=Stealth]
\def\surfacepath{(-2.4,-1.2) .. controls (-2.7,0.5) and (-1.9,2.2) .. (0,2.4)
  .. controls (1.9,2.6) and (2.9,1.2) .. (2.7,-0.2)
  .. controls (2.5,-1.5) and (0.7,-2.0) .. (-0.5,-1.8)
  .. controls (-1.6,-1.6) and (-2.1,-1.7) .. (-2.4,-1.2)}
\newcommand{\flen}{0.5}
% Gauge 1
\begin{scope}[shift={(-4.6,0)}]
  \node[above, font=\bfseries, text=gdlblue!70!black] at (0,2.9) {Gauge 1};
  \fill[gdlblue!15] \surfacepath; \draw[thick, gdlblue!40] \surfacepath;
  \foreach \p in {(-1.2,1.0),(1.0,1.2),(-0.9,-0.6),(1.1,-0.4),(0.0,0.3)} {
    \draw[->, thick, gdlgreen!50!black] \p -- ++(0:\flen);
    \draw[->, thick, gdlgreen!50!black] \p -- ++(90:\flen);
    \fill[gdlred!70] \p circle (2pt);
  }
\end{scope}
% Gauge 2
\begin{scope}[shift={(4.6,0)}]
  \node[above, font=\bfseries, text=gdlorange!70!black] at (0,2.9) {Gauge 2};
  \fill[gdlblue!15] \surfacepath; \draw[thick, gdlblue!40] \surfacepath;
  \draw[->, thick, gdlorange!70] (-1.2,1.0) -- ++(35:\flen);
  \draw[->, thick, gdlorange!70] (-1.2,1.0) -- ++(125:\flen);
  \draw[->, thick, gdlorange!70] (1.0,1.2) -- ++(200:\flen);
  \draw[->, thick, gdlorange!70] (1.0,1.2) -- ++(290:\flen);
  \draw[->, thick, gdlorange!70] (-0.9,-0.6) -- ++(110:\flen);
  \draw[->, thick, gdlorange!70] (-0.9,-0.6) -- ++(200:\flen);
  \draw[->, thick, gdlorange!70] (1.1,-0.4) -- ++(250:\flen);
  \draw[->, thick, gdlorange!70] (1.1,-0.4) -- ++(340:\flen);
  \draw[->, thick, gdlorange!70] (0.0,0.3) -- ++(60:\flen);
  \draw[->, thick, gdlorange!70] (0.0,0.3) -- ++(150:\flen);
  \foreach \p in {(-1.2,1.0),(1.0,1.2),(-0.9,-0.6),(1.1,-0.4),(0.0,0.3)}
    \fill[gdlred!70] \p circle (2pt);
\end{scope}
% arrow between
\draw[<->, very thick, gdlpurple!70] (-1.5,0) -- node[above, font=\small] {$g(x)\in\SO(2)$} (1.5,0);
\end{tikzpicture}
\caption[Libertad gauge]{Libertad gauge sobre una superficie. La misma superficie
admite distintas elecciones de marco local (Gauge 1, con orientaci\'on
consistente, y Gauge 2, rotado de forma independiente en cada punto). Una
convoluci\'on gauge-equivariante produce el mismo resultado geom\'etrico para
cualquiera de los dos, estando ambos relacionados por una transformaci\'on punto
a punto $g(x) \in \SO(2)$.}
\label{fig:gauge:freedom}
\end{figure}

Las caracter\'isticas en la variedad no son n\'umeros desnudos sino componentes
relativas al marco elegido, de modo que deben transformarse cuando lo hace el
marco. Esto se captura mediante el tipo de un campo de caracter\'isticas, fijado
por una representaci\'on $\rho$ del grupo de estructura.

\begin{definition}[Campo de caracter\'isticas gauge-equivariante]
\label{def:gauge:field}
Sea $\rho\colon G \to \GL(V)$ una representaci\'on. Un campo $f$ de \emph{tipo}
$\rho$ asigna a cada punto componentes en $V$ que, bajo una transformaci\'on
gauge $g$, se transforman como
\[
  f(p) \longmapsto \rho\big(g(p)^{-1}\big)\, f(p).
\]
Equivalentemente, considerando $f$ como una funci\'on en el espacio total,
$f(u \cdot g) = \rho(g^{-1})\, f(u)$ para todo $u \in P$ y $g \in G$.
\end{definition}

Casos especiales importantes en una superficie, con $G = \SO(2)$, son los campos
\emph{escalares} ($\rho$ trivial, invariantes bajo gauge), los campos
\emph{vectoriales} (la representaci\'on est\'andar $2$-dimensional, con
componentes que rotan con el marco), y los campos complejos que se transforman
por $e^{i k\theta}$ para un ``esp\'in'' entero $k$. Para un campo vectorial con
componentes $\mathbf{u} = u_1 e_1^v + u_2 e_2^v$, una rotaci\'on del marco por
$\theta_v$ act\'ua como
\[
  \begin{pmatrix} u_1' \\ u_2' \end{pmatrix} =
  \begin{pmatrix} \cos\theta_v & \sin\theta_v \\ -\sin\theta_v & \cos\theta_v
  \end{pmatrix}
  \begin{pmatrix} u_1 \\ u_2 \end{pmatrix},
\]
la ley de transformaci\'on est\'andar de un campo de tipo $1$.

El objetivo de dise\~no se sigue de inmediato. Una capa de red es
\emph{gauge-equivariante} si conmuta con la acci\'on gauge: transformar los
marcos de entrada produce la salida transformada de forma correspondiente. Si la
lectura final es una predicci\'on escalar global, la red es \emph{gauge-invariante},
\[
  \Phi(g \cdot f) = \Phi(f) \qquad \text{para toda transformaci\'on gauge } g.
\]
Esto garantiza que las predicciones son independientes de la parametrizaci\'on o
carta local, que las salidas con valores vectoriales rotan coherentemente con el
marco y --- en el caso especial del gauge de permutaci\'on de nodos en grafos ---
que las salidas son invariantes bajo el reetiquetado de los nodos.
\Cref{fig:gauge:transformation} muestra una transformaci\'on gauge actuando sobre
una se\~nal de grafo. La conexi\'on con la f\'isica es exacta: las teor\'ias gauge
describen las fuerzas fundamentales mediante conexiones en fibrados principales
con grupos de estructura $\Ut(1)$, $\SU(2)$, $\SU(3)$; en el aprendizaje profundo
geom\'etrico el grupo de estructura codifica las simetr\'ias locales relevantes
para el problema de aprendizaje.

\begin{figure}[ht]
\centering
\begin{tikzpicture}[>=Stealth,
  vertex/.style={circle, draw, thick, minimum size=22pt}]
% original
\begin{scope}[shift={(-4.6,0)}]
  \node[above, font=\bfseries] at (1,2.4) {Caracter\'isticas en un grafo};
  \node[vertex, fill=gdlblue!20] (a) at (0,0) {$v_1$};
  \node[vertex, fill=gdlblue!20] (b) at (2,0) {$v_2$};
  \node[vertex, fill=gdlblue!20] (c) at (1,1.5) {$v_3$};
  \draw[thick] (a)--(b); \draw[thick] (b)--(c); \draw[thick] (c)--(a);
  \draw[->, thick, gdlred!70] (a) -- ++(-0.5,0.3) node[left, font=\tiny] {$f_1$};
  \draw[->, thick, gdlgreen!60!black] (b) -- ++(0.5,0.2) node[right, font=\tiny] {$f_2$};
  \draw[->, thick, gdlorange!70] (c) -- ++(0.2,0.5) node[above, font=\tiny] {$f_3$};
\end{scope}
\draw[->, very thick, gdlpurple!70] (-2,0.6) -- node[above, font=\small] {gauge $g:V\to G$} (2,0.6);
% transformed
\begin{scope}[shift={(4.6,0)}]
  \node[above, font=\bfseries] at (1,2.4) {Caracter\'isticas transformadas};
  \node[vertex, fill=gdlblue!20] (a2) at (0,0) {$v_1$};
  \node[vertex, fill=gdlblue!20] (b2) at (2,0) {$v_2$};
  \node[vertex, fill=gdlblue!20] (c2) at (1,1.5) {$v_3$};
  \draw[thick] (a2)--(b2); \draw[thick] (b2)--(c2); \draw[thick] (c2)--(a2);
  \draw[->, thick, gdlred!70] (a2) -- ++(0.3,0.5) node[left, font=\tiny] {$\rho(g_1)f_1$};
  \draw[->, thick, gdlgreen!60!black] (b2) -- ++(-0.2,0.5) node[right, font=\tiny] {$\rho(g_2)f_2$};
  \draw[->, thick, gdlorange!70] (c2) -- ++(-0.4,0.3) node[above, font=\tiny] {$\rho(g_3)f_3$};
  \draw[dashed, gdlpurple!50] (a2) ++(0.3,0) arc (0:60:0.3cm);
  \draw[dashed, gdlpurple!50] (b2) ++(0.3,0) arc (0:120:0.3cm);
  \draw[dashed, gdlpurple!50] (c2) ++(0.3,0) arc (0:150:0.3cm);
\end{scope}
\end{tikzpicture}
\caption[Transformaci\'on gauge en un grafo]{Una transformaci\'on gauge $g\colon V
\to G$ rota el marco de referencia de forma independiente en cada v\'ertice,
enviando la caracter\'istica $f_i$ a $\rho(g_i)\, f_i$. Los arcos discontinuos
indican la rotaci\'on local aplicada en cada nodo.}
\label{fig:gauge:transformation}
\end{figure}

\begin{example}[Por qu\'e importa la equivarianza gauge]
\label{ex:gauge:mesh}
En una malla de tri\'angulos $3$D $\Mani = (V, F)$ el plano tangente $T_v\Mani$ en
cada v\'ertice $v$ est\'a bien definido, pero no hay base ortonormal can\'onica
$(e_1^v, e_2^v)$ para \'el: cualquier rotaci\'on $R_\theta \in \SO(2)$ da un marco
igualmente v\'alido. Un perceptr\'on multicapa ingenuo que prediga, digamos, la
direcci\'on de m\'axima curvatura a partir de las componentes $(u_1, u_2)$ de
vectores en coordenadas locales arbitrarias falla en silencio. Supongamos que con
el marco original predice la direcci\'on $(\hat d_1, \hat d_2) = (0.8, 0.6)$;
despu\'es de rotar ese marco por $45^\circ$ --- una transformaci\'on gauge ---
podr\'ia predecir $(0.3, 0.9)$, una direcci\'on completamente distinta. Una red
gauge-equivariante como la Gauge Equivariant Mesh CNN de
\citet{cohen2019gaugeequivariantconvolutional} garantiza en cambio
$\hat{\mathbf{d}}' = R_\theta\,\hat{\mathbf{d}}$ bajo una rotaci\'on del marco
$R_\theta$: las componentes num\'ericas cambian, pero la direcci\'on geom\'etrica
predicha no.
\end{example}

\begin{example}[Campos tensoriales en mol\'eculas $3$D]
\label{ex:gauge:molecules}
Al modelar una mol\'ecula con \'atomos en posiciones $\{\mathbf{r}_i\}_{i=1}^n
\subset \mathbb{R}^3$ se predicen a menudo cantidades tensoriales, como el tensor
de polarizabilidad $\boldsymbol{\alpha} \in \mathbb{R}^{3\times 3}$. Bajo una
rotaci\'on global $R \in \SO(3)$ este tensor se transforma como
$\boldsymbol{\alpha}' = R\, \boldsymbol{\alpha}\, R^{\!\top}$, una representaci\'on
tensorial de rango $2$. Redes equivariantes como NequIP~\citep{batzner2022nequip}
representan las caracter\'isticas intermedias como \emph{tensores esf\'ericos} que
se transforman bajo representaciones irreducibles de $\SO(3)$ y construyen las
salidas mediante productos tensoriales equivariantes,
\[
  \boldsymbol{\alpha} = \sum_{\ell = 0, 2} W_\ell \big(\mathbf{Y}^{(\ell)} \otimes
  \mathbf{Y}^{(\ell)}\big),
\]
con $\mathbf{Y}^{(\ell)}$ los arm\'onicos esf\'ericos de orden $\ell$ y $W_\ell$
pesos aprendidos. Esto garantiza autom\'aticamente el comportamiento rotacional
correcto, eliminando la necesidad de aumento de datos rotacional y produciendo
predicciones f\'isicamente consistentes e independientes de la orientaci\'on
arbitraria de la mol\'ecula.
\end{example}

\section{Curvatura, holonom\'ia y obstrucciones topol\'ogicas}
\label{sec:gauge:curvature}

El transporte paralelo a lo largo de un camino nos permite comparar marcos en los
extremos, pero el resultado puede depender del camino tomado. La
\emph{curvatura} de una conexi\'on mide esta dependencia del camino,
equivalentemente el fallo de la distribuci\'on horizontal en ser integrable.

\begin{definition}[Curvatura y la ecuaci\'on de estructura de Cartan]
\label{def:gauge:curvature}
La \emph{curvatura} de una conexi\'on $\omega$ es la $2$-forma con valores en
$\Lie{g}$
\[
  \Omega = \mathrm{d}\omega + \tfrac{1}{2}[\omega \wedge \omega],
\]
donde $[\omega \wedge \omega](X, Y) = [\omega(X), \omega(Y)]$ usa el corchete de
Lie; expl\'icitamente $\Omega(X, Y) = \mathrm{d}\omega(X, Y) + [\omega(X),
\omega(Y)]$. Es una $2$-forma horizontal y equivariante. Reordenada, la
definici\'on es la \emph{ecuaci\'on de estructura de Cartan}
\[
  \mathrm{d}\omega = \Omega - \tfrac{1}{2}[\omega \wedge \omega].
\]
\end{definition}

\begin{theorem}[Identidad de Bianchi]
\label{thm:gauge:bianchi}
La curvatura satisface $\mathrm{d}\Omega = [\omega \wedge \Omega]$.
\end{theorem}

\begin{proof}
Aplicando $\mathrm{d}$ a la definici\'on de $\Omega$ y usando $\mathrm{d}^2 = 0$,
\[
  \mathrm{d}\Omega = \mathrm{d}(\mathrm{d}\omega) + \tfrac{1}{2}\mathrm{d}[\omega
  \wedge \omega] = \tfrac{1}{2}\big([\mathrm{d}\omega \wedge \omega] - [\omega
  \wedge \mathrm{d}\omega]\big).
\]
Sustituyendo $\mathrm{d}\omega = \Omega - \tfrac{1}{2}[\omega \wedge \omega]$ y
usando la identidad de Jacobi para el corchete de Lie, los t\'erminos
cuadr\'aticos se cancelan, dejando $\mathrm{d}\Omega = [\omega \wedge \Omega]$.
\end{proof}

\begin{definition}[Holonom\'ia]
\label{def:gauge:holonomy}
Para un lazo $\gamma$ basado en $p$ (de modo que $\gamma(0) = \gamma(1) = p$), la
\emph{holonom\'ia} es el elemento $g_\gamma \in G$ con $\Gamma_\gamma(u) = u \cdot
g_\gamma$ para todo $u \in \pi^{-1}(p)$. El \emph{grupo de holonom\'ia}
$\Hol_p(\omega)$ es el subgrupo de $G$ generado por las holonom\'ias de todos los
lazos basados en $p$.
\end{definition}

La curvatura es la holonom\'ia en el l\'imite infinitesimal: para un lazo
peque\~no que encierra un \'area $\mathrm{d}A$ la holonom\'ia es aproximadamente
$\exp(\Omega \cdot \mathrm{d}A)$, y el teorema de Ambrose--Singer lo precisa
identificando $\Hol_p(\omega)$ con el subgrupo generado por los valores de la
curvatura $\Omega(X, Y)$. Cuando la curvatura se anula la conexi\'on es
\emph{plana} y el transporte paralelo depende solo de la clase de homotop\'ia del
camino; un vector transportado alrededor de un lazo contr\'actil regresa sin
cambios. En un dominio curvo la curvatura es generalmente no nula, y esto es
exactamente por lo que un marco transportado alrededor de un lazo cerrado regresa
rotado --- la obstrucci\'on que una convoluci\'on gauge-equivariante debe
respetar. \Cref{fig:gauge:holonomy} contrasta los casos plano y curvo.

\begin{figure}[ht]
\centering
\begin{tikzpicture}[>=Stealth, font=\small]
% flat
\begin{scope}[shift={(-4.2,0)}]
  \node[above, font=\bfseries] at (0,2.6) {Conexi\'on plana ($\Omega=0$)};
  \fill[gdlblue!18] (-1.6,-1.6) rectangle (1.6,1.6);
  \draw[thick, gdlblue!30] (-1.6,-1.6) rectangle (1.6,1.6);
  \draw[very thick, gdlred!70] (0,0) -- (1.2,0) -- (1.2,1.2) -- (0,1.2) -- cycle;
  \draw[->, thick, gdlgreen!60!black] (0,0) -- (0.5,0.25);
  \draw[->, thick, gdlgreen!60!black] (1.2,0) -- (1.7,0.25);
  \draw[->, thick, gdlgreen!60!black] (1.2,1.2) -- (1.7,1.45);
  \draw[->, thick, gdlgreen!60!black] (0,1.2) -- (0.5,1.45);
  \node[font=\footnotesize] at (0,-2.0) {$v_{\text{final}} = v_{\text{inicial}}$};
\end{scope}
% curved
\begin{scope}[shift={(4.2,0)}]
  \node[above, font=\bfseries] at (0,2.6) {Conexi\'on curva ($\Omega\neq0$)};
  \shade[ball color=gdlpurple!20, opacity=0.8] (0,0) circle (1.6cm);
  \draw[thick, gdlpurple!40] (0,0) circle (1.6cm);
  \draw[very thick, gdlred!70] (0,1.6) arc (90:0:1.6cm) arc (0:-90:1.6cm and 0.4cm) arc (-90:90:0.4cm and 1.6cm);
  \draw[->, thick, gdlgreen!60!black] (0,1.6) -- (0.45,1.85);
  \draw[->, thick, gdlgreen!60!black] (1.6,0) -- (1.9,0.45);
  \draw[->, thick, gdlgreen!60!black] (0,-0.4) -- (0.45,-0.15);
  \draw[->, very thick, gdlorange!70, dashed] (0,1.6) -- (-0.3,1.95);
  \draw[thick, gdlorange!70] (0.25,1.7) arc (30:120:0.3cm);
  \node[gdlorange!70, font=\tiny] at (-0.1,2.05) {$\theta$};
  \node[font=\footnotesize] at (0,-2.0) {$v_{\text{final}} \neq v_{\text{inicial}}$};
\end{scope}
\node[font=\small, text width=8cm, align=center] at (0,-3.0)
  {\textbf{Ambrose--Singer:}\ $\Hol_p(\omega) = \big\langle \Omega(X,Y) : X,Y \in
  T_u P\big\rangle$ --- la holonom\'ia est\'a generada por la curvatura.};
\end{tikzpicture}
\caption[Holonom\'ia y curvatura]{Transporte paralelo alrededor de un lazo
cerrado. Para una conexi\'on plana (izquierda) el vector transportado regresa sin
cambios y la holonom\'ia es trivial. Para una conexi\'on curva (derecha) regresa
rotado por un \'angulo $\theta$; la holonom\'ia equivale a la integral de la
curvatura sobre la superficie encerrada.}
\label{fig:gauge:holonomy}
\end{figure}

La curvatura tambi\'en explica \emph{por qu\'e} la equivarianza gauge es necesaria
en lugar de meramente conveniente.

\begin{proposition}[Obstrucci\'on topol\'ogica a un gauge global]
\label{prop:gauge:obstruction}
Un fibrado principal $P \to \Mani$ admite una secci\'on global --- una elecci\'on
de gauge consistente sobre todo $\Mani$ --- si y solo si es trivial,
$P \cong \Mani \times G$. En particular:
\begin{enumerate}
  \item el fibrado de marcos ortonormales de $S^2$ no admite secci\'on global
    (por el teorema de la bola peluda);
  \item todo fibrado principal sobre una variedad contr\'actil es trivial.
\end{enumerate}
\end{proposition}

\begin{proof}
Una secci\'on global $s\colon \Mani \to P$ induce la trivializaci\'on
$\Mani \times G \to P$, $(p, g) \mapsto s(p) \cdot g$, un isomorfismo de fibrados
por la acci\'on libre y transitiva de $G$ sobre cada fibra; rec\'iprocamente una
trivializaci\'on $P \cong \Mani \times G$ proporciona la secci\'on $s(p) = (p,
e)$. Para (1), una secci\'on del fibrado de marcos suministrar\'ia un campo
continuo de marcos ortonormales, en particular un campo vectorial tangente que no
se anula en ning\'un punto, contradiciendo el teorema de la bola peluda
($\chi(S^2) = 2 \neq 0$). Para (2), una base contr\'actil es homot\'opicamente
equivalente a un punto, y los fibrados principales se clasifican por clases de
homotop\'ia de aplicaciones al espacio clasificante $BG$, de modo que todo
fibrado sobre ella es trivial.
\end{proof}

La consecuencia para el aprendizaje es decisiva: en $S^2$ y en la mayor\'ia de
las variedades de inter\'es no hay un sistema de coordenadas globalmente
consistente, y cualquier parametrizaci\'on global tiene singularidades (los polos
en coordenadas esf\'ericas). La equivarianza gauge libera a la red de estas
elecciones locales arbitrarias y evita los artefactos en las singularidades, en
lugar de disimularlas con una carta fija.

\section{Convoluci\'on gauge-equivariante}
\label{sec:gauge:conv}

Ahora ensamblamos las piezas en una convoluci\'on. La dificultad que una capa
gauge-equivariante debe superar es concreta: para combinar la caracter\'istica en
un punto $p$ con las caracter\'isticas de sus vecinos, las caracter\'isticas de
esos vecinos viven en fibras distintas, expresadas en marcos locales distintos, y
no pueden sumarse directamente. La conexi\'on resuelve esto --- el transporte
paralelo lleva el marco de cada vecino a $p$ antes de combinar las
caracter\'isticas. Primero formalizamos la distinci\'on entre transformaciones
globales y locales que subyace a la construcci\'on.

\begin{definition}[Transformaci\'on local frente a global]
\label{def:gauge:local-global}
Para un grupo $G$ y un campo $f\colon \Mani \to \mathbb{R}^d$:
una transformaci\'on \emph{global} es $T_g[f](x) = \rho(g)\, f(g^{-1} \cdot x)$,
con un \'unico $g \in G$ usado en cada punto $x$; una transformaci\'on
\emph{local} (gauge) es $T_{g(\cdot)}[f](x) = \rho(g(x))\, f(\phi_{g(x)}(x))$, con
$g\colon \Mani \to G$ asignando un elemento posiblemente distinto a cada punto.
\end{definition}

\begin{theorem}[Convoluci\'on gauge-equivariante]
\label{thm:gauge:conv}
Sea $f$ un campo de caracter\'isticas en $\Mani$ y sea $\Gamma_{q \to p}$ el
transporte paralelo de $q$ a $p$. La convoluci\'on
\[
  (f \star \psi)(p) = \int_{\mathcal{N}(p)}
    \psi\big(q, p\big)\; \Gamma_{q \to p}\, f(q)\; \mathrm{d}q,
\]
que transporta la caracter\'istica de cada vecino al marco en $p$ antes de
aplicar el kernel aprendible $\psi$, mapea campos gauge-equivariantes a campos
gauge-equivariantes.
\end{theorem}

\begin{proof}
Sea $\rho_p$ un cambio de gauge en $p$, de modo que los campos de
caracter\'isticas se transforman como $f(q) \mapsto \rho_q\, f(q)$. El operador de
transporte $\Gamma_{q \to p}$ cambia de base del marco en $q$ al marco en $p$,
luego se transforma como $\Gamma_{q \to p} \mapsto \rho_p\, \Gamma_{q \to p}\,
\rho_q^{-1}$. Sustituyendo,
\[
  (f \star \psi)(p) \mapsto \int_{\mathcal{N}(p)} \psi(q,p)\,
    \big(\rho_p\, \Gamma_{q\to p}\, \rho_q^{-1}\big)\,\big(\rho_q\, f(q)\big)\,
    \mathrm{d}q
  = \rho_p \int_{\mathcal{N}(p)} \psi(q,p)\, \Gamma_{q\to p}\, f(q)\, \mathrm{d}q,
\]
ya que $\rho_q^{-1}\rho_q = \mathrm{Id}$. As\'i $(f \star \psi)(p) \mapsto \rho_p\,
(f \star \psi)(p)$, la ley de transformaci\'on de un campo
equivariante~\citep{cohen2019gaugeequivariantconvolutional}.
\end{proof}

La cancelaci\'on es el mismo mecanismo algebraico que hace que la intensidad de
campo $F_{\mu\nu} = \partial_\mu A_\nu - \partial_\nu A_\mu$ sea gauge-invariante
en electromagnetismo aunque el potencial $A_\mu$ no lo sea: los factores gauge en
el punto fuente se cancelan contra los factores inversos transportados por la
conexi\'on. De forma crucial, la equivarianza gauge impone restricciones
\emph{lineales} sobre el kernel $\psi$, reduciendo el espacio de par\'ametros a
un subespacio af\'in que puede resolverse anal\'iticamente antes del
entrenamiento --- el mismo principio de restricci\'on de kernel que hace
pr\'acticas las convoluciones de grupo orientables (steerable) en
\Cref{ch:groups}.

\subsection{La CNN icosa\'edrica}
\label{sec:gauge:icosahedral}

La brecha entre la teor\'ia suave y una red finita y entrenable fue salvada por
primera vez en la pr\'actica por \citet{cohen2019gaugeequivariantconvolutional},
quien eligi\'o el icosaedro como discretizaci\'on de la esfera $S^2$. La
elecci\'on dista de ser arbitraria: entre los poliedros regulares el icosaedro es
el que mejor aproxima la esfera, tiene un grupo de simetr\'ia rico con $120$
elementos, sus v\'ertices est\'an espaciados de forma casi uniforme, y con solo
$12$ v\'ertices, $20$ caras y $30$ aristas (refinables por subdivisi\'on
geod\'esica) es computacionalmente ligero.

En esta malla el grupo de estructura es $\SO(2)$ --- rotaciones planas del marco
tangente en cada v\'ertice. Un campo de caracter\'isticas discreto $f\colon V \to
V_\rho$ vive en los v\'ertices, y la convoluci\'on gauge-equivariante se convierte
en una suma local sobre los vecinos,
\[
  (f \star \psi)(v_i) = \sum_{j \in \mathcal{N}(i)}
    \psi_{ij}\; \Gamma_{ij}\, f(v_j),
\]
donde $\Gamma_{ij} \in \SO(2)$ es el transporte paralelo discreto del v\'ertice
$v_j$ al $v_i$ a lo largo de la arista que los conecta --- una rotaci\'on por un
\'angulo fijo determinado por la geometr\'ia de la malla --- y $\psi_{ij}$ son los
pesos aprendibles del kernel, restringidos para que la capa respete la acci\'on
gauge de $\SO(2)$. El transporte de una caracter\'istica desde un v\'ertice
distante se obtiene componiendo las rotaciones por arista $\Gamma_{ij} =
\exp(-\omega_e)$ a lo largo de un camino geod\'esico, exactamente la contraparte
discreta de la exponencial ordenada por camino de \Cref{sec:gauge:transport}:
inicializar el transporte a la identidad y multiplicar por $\exp(-\omega_e)$ por
cada arista $e$ recorrida. Para evitar soluciones degeneradas se penaliza el
exceso de curvatura mediante un regularizador que suma los \'angulos de las
aristas alrededor de cada tri\'angulo,
\[
  \mathcal{R}(\theta) = \sum_{\text{tri\'angulos } T}
    \Big|\sum_{(i,j) \in \partial T} \theta_{ij}\Big|^2,
\]
una lectura discreta de la identidad de Bianchi de \Cref{thm:gauge:bianchi}:
controla la holonom\'ia acumulada alrededor de lazos peque\~nos.

\Cref{fig:gauge:icosahedral} esboza el mecanismo. El beneficio pr\'actico
decisivo es que todo este c\'omputo puede expresarse en t\'erminos de operaciones
est\'andar tipo convoluci\'on sobre una representaci\'on aplanada del icosaedro,
haciendo la red tan eficiente como una CNN ordinaria a la vez que es
\emph{exactamente} gauge-equivariante --- y, como el grupo de simetr\'ia
icosa\'edrico act\'ua sobre la malla, equivariante tambi\'en a ese subgrupo
discreto de rotaciones de la esfera.

\begin{figure}[ht]
\centering
\begin{tikzpicture}[>=Stealth, font=\small]
% central vertex and neighbours (hexagonal local patch)
\coordinate (c) at (0,0);
\foreach \a/\name in {0/n0,60/n1,120/n2,180/n3,240/n4,300/n5}
  \coordinate (\name) at (\a:1.8);
% edges
\foreach \name in {n0,n1,n2,n3,n4,n5} \draw[gdlgray!55, thick] (c) -- (\name);
\foreach \a/\b in {n0/n1,n1/n2,n2/n3,n3/n4,n4/n5,n5/n0}
  \draw[gdlgray!35] (\a) -- (\b);
% transported frames (rotated) at neighbours
\foreach \name/\rot in {n0/20,n1/-15,n2/40,n3/10,n4/-30,n5/25} {
  \draw[->, thick, gdlorange!75] (\name) -- ++(\rot:0.6);
  \draw[->, thick, gdlorange!75] (\name) -- ++({\rot+90}:0.6);
  \fill[gdlblue!65] (\name) circle (2.5pt);
}
% transport arrows toward centre
\foreach \name in {n0,n2,n4}
  \draw[->, gdlpurple!70, thick, dashed] ($(\name)!0.30!(c)$) -- ($(\name)!0.62!(c)$);
% central frame
\draw[->, very thick, gdlgreen!55!black] (c) -- ++(0:0.7);
\draw[->, very thick, gdlgreen!55!black] (c) -- ++(90:0.7);
\fill[gdlred!75] (c) circle (3pt);
\node[below right, font=\footnotesize, text=gdlgreen!45!black] at (c) {$v_i$};
\node[font=\footnotesize, text=gdlpurple!70] at (1.15,0.55) {$\Gamma_{ij}$};
\node[font=\footnotesize, text=gdlorange!70!black] at (-1.9,1.55) {$f(v_j)$};
\end{tikzpicture}
\caption[Convoluci\'on gauge-equivariante en una malla]{Convoluci\'on
gauge-equivariante local en un v\'ertice $v_i$. Cada vecino lleva una
caracter\'istica en su propio marco local (naranja). Antes de combinar, la
conexi\'on la transporta paralelamente al marco en $v_i$ mediante la rotaci\'on
$\Gamma_{ij} \in \SO(2)$ (p\'urpura), tras lo cual se aplica el kernel aprendible
$\psi_{ij}$. El resultado se transforma coherentemente bajo cualquier cambio de
marco local.}
\label{fig:gauge:icosahedral}
\end{figure}

\subsection{El marco general y perspectivas}
\label{sec:gauge:outlook}

El caso icosa\'edrico instancia una receta general. Elegir una variedad base
$\Mani$ para los datos y un grupo de estructura $G$ para las simetr\'ias locales
relevantes fija un fibrado principal; una conexi\'on suministra el transporte
paralelo; y la convoluci\'on gauge-equivariante de \Cref{thm:gauge:conv} produce
una capa cuyo tipo de salida est\'a gobernado por una representaci\'on elegida
$\rho$. Distintos problemas requieren distintos fibrados: fibrados triviales
cuando las simetr\'ias son globalmente consistentes, fibrados de marcos cuando se
requiere invarianza bajo cambios locales de coordenadas, fibrados de esp\'in para
datos con estructura de orientaci\'on, y fibrados gauge m\'as ricos para campos
con simetr\'ias internas (electromagnetismo con $\Ut(1)$, teor\'ias de
Yang--Mills con grupos no abelianos, cromodin\'amica cu\'antica en ret\'iculo con
$\SU(3)$). Bajo condiciones suaves tales redes son aproximadores universales de
aplicaciones gauge-equivariantes: el espacio de funciones realizables es denso en
el espacio $L^2$ de campos gauge-equivariantes, con tasa de aproximaci\'on
$O(n^{-2/(d_\Mani + d_G)})$ gobernada por las dimensiones de la base y del grupo y
constantes que dependen de la curvatura del fibrado.

Este marco es la culminaci\'on de las cinco G. Contiene estrictamente a las
dem\'as: un grupo de simetr\'ia global es el caso especial de un fibrado trivial
con una conexi\'on plana, recuperando las redes equivariantes a grupos de
\Cref{ch:groups}; una variedad con su conexi\'on de Levi-Civita recupera las
convoluciones geod\'esicas de \Cref{ch:geodesics} como una elecci\'on de gauge
particular; y las redes esf\'ericas que motivaron este desarrollo se sit\'uan
dentro de \'el como la instancia icosa\'edrica. El anidamiento
\[
  \mathcal{F}_{\text{eucl\'ideo}} \subset \mathcal{F}_{\text{grupo}}
  \subset \mathcal{F}_{\text{geod\'esico}} \subset \mathcal{F}_{\text{gauge}}
\]
resume c\'omo cada una de las cuatro G anteriores es una versi\'on restringida de
la convoluci\'on gauge-equivariante.

\section{\'Algebras de Clifford y geom\'etricas}
\label{sec:gauge:clifford}

La perspectiva gauge admite una potente generalizaci\'on algebraica. Donde un
campo de caracter\'isticas expresa una cantidad direccional en componentes
relativas a un marco, un \emph{\'algebra geom\'etrica (de Clifford)} permite que
escalares, vectores, \'areas orientadas y vol\'umenes orientados vivan en un
\'unico objeto graduado y se transformen juntos bajo rotaciones y reflexiones.
Esto unifica los n\'umeros complejos, los cuaterniones y las representaciones
tensoriales, y proporciona bloques de construcci\'on equivariantes tanto para la
convoluci\'on como para la atenci\'on.

\begin{definition}[\'Algebra de Clifford]
\label{def:gauge:clifford}
Sea $V$ un espacio vectorial real de dimensi\'on finita con forma cuadr\'atica
$Q$. El \emph{\'algebra de Clifford} $\Cl(V, Q)$ es el \'algebra asociativa
unitaria generada por $V$ sujeta a $v^2 = Q(v)\, 1$ para todo $v \in V$. Para una
m\'etrica de signatura $(p, q)$ escribimos $\Cl(p, q)$, con $p$ direcciones de
norma positiva y $q$ de norma negativa.
\end{definition}

El \'algebra unifica varias estructuras familiares: $\Cl(0,0) \cong \mathbb{R}$,
$\Cl(0,1) \cong \mathbb{C}$, $\Cl(0,2) \cong \mathbb{H}$ (los cuaterniones), y
$\Cl(3,0)$ es el \'algebra geom\'etrica eucl\'idea de $\mathbb{R}^3$.

\begin{definition}[Multivector y graduaci\'on]
\label{def:gauge:multivector}
Los elementos de $\Cl(p, q)$ son \emph{multivectores} y se descomponen de forma
\'unica por grado. Con $n = p + q$,
\[
  M = \langle M \rangle_0 + \langle M \rangle_1 + \cdots + \langle M \rangle_n,
\]
donde $\langle M \rangle_0 \in \mathbb{R}$ es un escalar, $\langle M \rangle_1 \in
\operatorname{span}\{\mathbf{e}_1, \ldots, \mathbf{e}_n\}$ un vector,
$\langle M \rangle_2$ un bivector (\'area orientada), $\langle M \rangle_k$ un
$k$-vector de dimensi\'on $\binom{n}{k}$, y $\langle M \rangle_n$ el
pseudoescalar.
\end{definition}

\begin{theorem}[Dimensi\'on del \'algebra]
\label{thm:gauge:clifford-dim}
El \'algebra de Clifford $\Cl(p, q)$ es un espacio vectorial de dimensi\'on
$2^{p+q}$, con base can\'onica los productos ordenados de subconjuntos de
$\{\mathbf{e}_1, \ldots, \mathbf{e}_{p+q}\}$:
$\dim \Cl(p,q) = \sum_{k=0}^{p+q} \binom{p+q}{k} = 2^{p+q}$.
\end{theorem}

\begin{proof}
Sea $n = p+q$ y $\{\mathbf{e}_i\}$ una base ortonormal. Los multivectores base
son los productos ordenados $\mathbf{e}_{i_1}\cdots\mathbf{e}_{i_k}$ con
$1 \le i_1 < \cdots < i_k \le n$, $0 \le k \le n$ (con $k=0$ el escalar $1$).
Son linealmente independientes: las relaciones $\mathbf{e}_i\mathbf{e}_j =
-\mathbf{e}_j\mathbf{e}_i$ para $i \neq j$ ordenan cualquier producto, y
$\mathbf{e}_i^2 = \pm 1$ elimina los cuadrados. El n\'umero de subconjuntos de
$\{1, \ldots, n\}$ es $\sum_{k=0}^{n}\binom{n}{k} = 2^n$, la dimensi\'on del
\'algebra.
\end{proof}

\begin{definition}[Producto geom\'etrico]
\label{def:gauge:geometric-product}
Para vectores $\mathbf{u}, \mathbf{v} \in V$ el \emph{producto geom\'etrico}
unifica los productos interior y exterior:
\[
  \mathbf{u}\mathbf{v} = \mathbf{u}\cdot\mathbf{v} + \mathbf{u}\wedge\mathbf{v},
  \qquad
  \mathbf{u}\cdot\mathbf{v} = \tfrac{1}{2}(\mathbf{u}\mathbf{v} +
  \mathbf{v}\mathbf{u}),
  \quad
  \mathbf{u}\wedge\mathbf{v} = \tfrac{1}{2}(\mathbf{u}\mathbf{v} -
  \mathbf{v}\mathbf{u}).
\]
\end{definition}

\begin{example}[Rotores y reflexiones en $\Cl(3,0)$]
\label{ex:gauge:rotors}
En $\Cl(2,0)$ una rotaci\'on por el \'angulo $\theta$ se codifica mediante el
\emph{rotor} $R = \cos(\theta/2) + \sin(\theta/2)\,\mathbf{e}_1\wedge\mathbf{e}_2 =
e^{(\theta/2)\mathbf{e}_1\wedge\mathbf{e}_2}$, que act\'ua sobre un vector por el
producto s\'andwich $\mathbf{v}' = R\,\mathbf{v}\,\widetilde{R}$, donde
$\widetilde{R}$ es el reverso. Una reflexi\'on a trav\'es del hiperplano
ortogonal a un vector unitario $\mathbf{n}$ es $\mathbf{v}' =
-\mathbf{n}\,\mathbf{v}\,\mathbf{n}$; por ejemplo $\mathbf{n} = \mathbf{e}_1$ y
$\mathbf{v} = a\mathbf{e}_1 + b\mathbf{e}_2$ dan $\mathbf{v}' = -a\mathbf{e}_1 +
b\mathbf{e}_2$, invirtiendo la componente perpendicular. Un bivector elevado al
cuadrado da $-1$ ($(\mathbf{e}_1\wedge\mathbf{e}_2)^2 = -1$), de modo que juega el
papel algebraico de la unidad imaginaria. Estas operaciones s\'andwich son
exactamente transformaciones gauge realizadas dentro del \'algebra, con los
rotores implementando $\SO(n)$ y las reflexiones generando el $\Ot(n)$ completo.
\end{example}

\begin{definition}[Convoluci\'on de Clifford y fibrado de Clifford]
\label{def:gauge:clifford-conv}
Para una se\~nal con valores en Clifford $f\colon \Omega \to \Cl(p,q)$ y un kernel
$\psi\colon G \to \Cl(p,q)$ sobre un grupo de transformaciones $G$, la
\emph{convoluci\'on de Clifford} es
\[
  (f \star \psi)(x) = \int_G \rho_g^{-1}\big(\psi(g)\big)\cdot f(g^{-1}\cdot x)\,
  \mathrm{d}g,
\]
con $\cdot$ el producto geom\'etrico y $\rho_g$ la acci\'on de $G$ sobre el
\'algebra. En una variedad riemanniana $(\Mani, g)$ el \emph{fibrado de Clifford}
$\Cl(T\Mani)$ tiene fibra $\Cl(T_p\Mani, g_p)$ en $p$, y los campos de
caracter\'isticas son sus secciones: $s(p) \in \Cl(T_p\Mani)$.
\end{definition}

Las convoluciones con valores en Clifford empaquetan varios aspectos geom\'etricos
en una sola operaci\'on: los grados de un multivector portan un significado
geom\'etrico distinto e interpretable; la equivarianza es autom\'atica a partir de
la estructura algebraica; y las operaciones se prestan a una implementaci\'on
vectorizada. Subsumen las convoluciones de grupo orientables (\Cref{ch:groups}):
cuando una representaci\'on $\rho$ se descompone en irreducibles isomorfos a
espacios de multivectores de grado fijo de $\Cl(n,0)$, las aplicaciones lineales
equivariantes se corresponden con operaciones de Clifford que preservan el grado,
recuperando la estructura en bloques de las CNN orientables. La elecci\'on del
\'algebra es en s\'i misma una decisi\'on de modelado: el \'algebra eucl\'idea
$\Cl(3,0)$ ($8$ dimensiones) es ligera pero limitada a $\SO(3)$; las \'algebras
proyectiva y conforme $\Cl(4,1)$ ($32$ dimensiones) representan traslaciones,
esferas e inversiones a mayor coste, y emp\'iricamente el \'algebra conforme da
los mejores resultados~\citep{dehaan2024euclidean}. Aprender la m\'etrica del
\'algebra de forma adaptativa (redes equivariantes al grupo de Clifford con
$G_\theta = \sum_i \lambda_i(\theta)\,\mathbf{v}_i\mathbf{v}_i^{\!\top}$ aprendible)
permite que la geometr\'ia del \'algebra se ajuste a los datos preservando la
equivarianza~\citep{ali2024metric}.

\paragraph{Geometric Algebra Transformers (GATr).}
La misma \'algebra impulsa la atenci\'on equivariante. GATr~\citep{brehmer2023geometric}
opera en el \'algebra geom\'etrica proyectiva $\Cl(3,0,1)$ (un \'algebra de
dimensi\'on $16$ sobre la base $\{\mathbf{e}_0, \mathbf{e}_1, \mathbf{e}_2,
\mathbf{e}_3\}$ con $\mathbf{e}_i^2 = +1$ para $i = 1,2,3$ y $\mathbf{e}_0^2 =
0$). La direcci\'on degenerada $\mathbf{e}_0$ es lo que hace expresable la
geometr\'ia af\'in --- incluyendo las traslaciones ---, no solo la geometr\'ia
lineal. Los objetos geom\'etricos de $\mathbb{R}^3$ reciben representaciones
multivectoriales can\'onicas: escalares en grado $0$, planos $\pi =
n_1\mathbf{e}_1 + n_2\mathbf{e}_2 + n_3\mathbf{e}_3 + d\,\mathbf{e}_0$ en grado
$1$, l\'ineas en grado $2$, puntos en grado $3$, y el pseudoescalar $I =
\mathbf{e}_0\wedge\mathbf{e}_1\wedge\mathbf{e}_2\wedge \mathbf{e}_3$ en grado $4$,
con dimensiones por grado $1, 4, 6, 4, 1$ que suman $16$. Las rotaciones y
traslaciones act\'uan mediante rotores a trav\'es del producto s\'andwich, de modo
que la atenci\'on construida a partir de aplicaciones lineales que preservan el
grado, productos geom\'etricos y un producto interior equivariante es exactamente
$\mathrm{E}(3)$-equivariante. GATr es la contraparte basada en mecanismo de
atenci\'on de la convoluci\'on gauge-equivariante: el grupo gauge se realiza
algebraicamente, mediante operaciones de Clifford, en lugar de mediante una
conexi\'on expl\'icita.

\section{Extensiones topol\'ogicas}
\label{sec:gauge:topology}

La descripci\'on por fibrados trata el dominio como una variedad suave, pero
muchos dominios son genuinamente combinatorios, con celdas de varias dimensiones
y relaciones de incidencia m\'as ricas que las aristas de un grafo. El
\emph{aprendizaje profundo topol\'ogico} extiende el paso de mensajes a tales
estructuras, y el \'algebra de los operadores de frontera desempe\~na el papel que
la conexi\'on desempe\~n\'o arriba: especifica c\'omo fluye la informaci\'on entre
celdas de dimensi\'on adyacente.

\subsection{Complejos celulares, cadenas y operadores de frontera}
\label{sec:gauge:chains}

Recordemos que un \emph{complejo simplicial} $\mathcal{K}$ es una familia cerrada
hacia abajo de conjuntos finitos que se intersecan por pares en caras comunes; un
$0$-s\'implice es un v\'ertice, un $1$-s\'implice una arista, un $2$-s\'implice un
tri\'angulo, un $3$-s\'implice un tetraedro. Un \emph{complejo celular (CW)}
generaliza esto pegando celdas $e^n$ (homeomorfas a bolas abiertas) a lo largo de
aplicaciones caracter\'isticas $\phi\colon \partial\mathbb{D}^n \to X^{n-1}$ al
esqueleto inferior; las celdas pueden tener forma arbitraria (un cubo es una sola
celda $3$ en vez de varios tetraedros). Un \emph{complejo combinatorio} relaja
a\'un m\'as a cualquier familia de celdas con una relaci\'on de inclusi\'on
parcial, subsumiendo grafos, hipergrafos, complejos simpliciales y
celulares~\citep{hajij2023topological}.

\begin{definition}[Cadenas y el operador de frontera]
\label{def:gauge:chains}
El espacio de \emph{$k$-cadenas} es $C_k(\mathcal{K}) = \{\sum_i a_i\sigma_i : a_i
\in \mathbb{R}, \sigma_i\ \text{un } k\text{-s\'implice}\}$. El \emph{operador de
frontera} $\partial_k\colon C_k \to C_{k-1}$ act\'ua sobre un $k$-s\'implice
$\sigma = [v_0, \ldots, v_k]$ por
\[
  \partial_k \sigma = \sum_{i=0}^k (-1)^i [v_0, \ldots, \hat v_i, \ldots, v_k],
\]
con $\hat v_i$ denotando omisi\'on, extendido linealmente. El operador de
\emph{cofrontera} es el adjunto $\delta_k = \partial_{k+1}^{\!\top}\colon C_k \to
C_{k+1}$.
\end{definition}

\begin{theorem}[$\partial^2 = 0$]
\label{thm:gauge:partial-squared}
Para todo $k$, $\partial_{k-1}\circ\partial_k = 0$: la frontera de una frontera es
vac\'ia.
\end{theorem}

\begin{proof}
Para $\sigma = [v_0, \ldots, v_k]$,
\[
  \partial_{k-1}(\partial_k\sigma) = \sum_{i} (-1)^i \Big(\sum_{j<i}(-1)^j
  [\ldots \hat v_j \ldots \hat v_i \ldots] + \sum_{j>i}(-1)^{j-1}
  [\ldots \hat v_i \ldots \hat v_j \ldots]\Big).
\]
Cada cara con dos v\'ertices omitidos aparece exactamente dos veces, una por cada
orden de eliminaci\'on, con signos opuestos, de modo que todos los t\'erminos se
cancelan.
\end{proof}

Por ejemplo $\partial_2[v_0, v_1, v_2] = [v_1, v_2] - [v_0, v_2] + [v_0, v_1]$,
las tres aristas orientadas del tri\'angulo, y aplicar $\partial_1$ da cero.
\Cref{fig:gauge:simplicial-hierarchy} muestra la jerarqu\'ia resultante de cadenas
enlazadas por aplicaciones de frontera y cofrontera.

\begin{figure}[ht]
\centering
\begin{tikzpicture}[>=Stealth, font=\small, scale=0.62, transform shape]
\def\leveltop{7}\def\levelmid{0}\def\levelbot{-7}
\begin{scope}[on background layer]
  \fill[gdlblue!6, rounded corners=8pt] (-2.6,\leveltop-2.2) rectangle (7.8,\leveltop+2.2);
  \fill[gdlorange!6, rounded corners=8pt] (-2.6,\levelmid-2.2) rectangle (7.8,\levelmid+2.2);
  \fill[gdlred!6, rounded corners=8pt] (-2.6,\levelbot-2.2) rectangle (7.8,\levelbot+2.2);
\end{scope}
% level 0: nodes
\begin{scope}[shift={(0,\leveltop)}]
  \node[font=\bfseries, text=gdlblue!70!black, anchor=west] at (-2.4,1.5) {$0$-cadenas (nodos), $L_0$};
  \foreach \a/\b in {0/0,1.8/0.6,3.2/-0.2,1.0/-1.5,2.8/-1.8,4.2/-1.2}
    \fill[gdlblue!70] (\a,\b) circle (4pt);
  \draw[gdlgray!40,thick] (0,0)--(1.8,0.6) (0,0)--(1.0,-1.5) (1.8,0.6)--(3.2,-0.2)
    (1.8,0.6)--(1.0,-1.5) (1.8,0.6)--(2.8,-1.8) (3.2,-0.2)--(2.8,-1.8)
    (3.2,-0.2)--(4.2,-1.2) (1.0,-1.5)--(2.8,-1.8) (2.8,-1.8)--(4.2,-1.2);
\end{scope}
% level 1: edges
\begin{scope}[shift={(0,\levelmid)}]
  \node[font=\bfseries, text=gdlorange!70!black, anchor=west] at (-2.4,1.5) {$1$-cadenas (aristas), $L_1$};
  \draw[gdlorange!70,very thick] (0,0)--(1.8,0.6) (0,0)--(1.0,-1.5) (1.8,0.6)--(3.2,-0.2)
    (1.8,0.6)--(1.0,-1.5) (1.8,0.6)--(2.8,-1.8) (3.2,-0.2)--(2.8,-1.8)
    (3.2,-0.2)--(4.2,-1.2) (1.0,-1.5)--(2.8,-1.8) (2.8,-1.8)--(4.2,-1.2);
  \foreach \a/\b in {0/0,1.8/0.6,3.2/-0.2,1.0/-1.5,2.8/-1.8,4.2/-1.2}
    \fill[gdlgray!50] (\a,\b) circle (3pt);
\end{scope}
% level 2: faces
\begin{scope}[shift={(0,\levelbot)}]
  \node[font=\bfseries, text=gdlred!70!black, anchor=west] at (-2.4,1.5) {$2$-cadenas (caras), $L_2$};
  \fill[gdlred!15] (0,0)--(1.8,0.6)--(1.0,-1.5)--cycle;
  \fill[gdlred!15] (1.8,0.6)--(1.0,-1.5)--(2.8,-1.8)--cycle;
  \fill[gdlred!15] (3.2,-0.2)--(2.8,-1.8)--(4.2,-1.2)--cycle;
  \draw[gdlred!60,thick] (0,0)--(1.8,0.6)--(1.0,-1.5)--cycle
    (1.8,0.6)--(1.0,-1.5)--(2.8,-1.8)--cycle (3.2,-0.2)--(2.8,-1.8)--(4.2,-1.2)--cycle;
  \foreach \a/\b in {0/0,1.8/0.6,3.2/-0.2,1.0/-1.5,2.8/-1.8,4.2/-1.2}
    \fill[gdlgray!50] (\a,\b) circle (3pt);
\end{scope}
% boundary / coboundary arrows
\draw[->, gdlpurple!70, line width=1.5pt] (-3.4,\levelmid+2.3) -- (-3.4,\leveltop-2.3);
\node[gdlpurple!70!black] at (-3.4,{(\levelmid+\leveltop)/2}) {$\partial_1$};
\draw[->, gdlgreen!70!black, line width=1.5pt] (-5.4,\leveltop-2.3) -- (-5.4,\levelmid+2.3);
\node[gdlgreen!70!black] at (-5.4,{(\leveltop+\levelmid)/2}) {$\delta_0$};
\draw[->, gdlpurple!70, line width=1.5pt] (-3.4,\levelbot+2.3) -- (-3.4,\levelmid-2.3);
\node[gdlpurple!70!black] at (-3.4,{(\levelbot+\levelmid)/2}) {$\partial_2$};
\draw[->, gdlgreen!70!black, line width=1.5pt] (-5.4,\levelmid-2.3) -- (-5.4,\levelbot+2.3);
\node[gdlgreen!70!black] at (-5.4,{(\levelmid+\levelbot)/2}) {$\delta_1$};
\end{tikzpicture}
\caption[Jerarqu\'ia simplicial]{Paso de mensajes multiorden en un complejo
simplicial: las se\~nales en nodos ($0$-cadenas), aristas ($1$-cadenas) y caras
($2$-cadenas) se enlazan mediante operadores de frontera $\partial_k$ y
operadores de cofrontera $\delta_k$. El laplaciano de Hodge $L_k$ acopla la
se\~nal de nivel $k$ con sus vecinos de dimensi\'on $k \pm 1$.}
\label{fig:gauge:simplicial-hierarchy}
\end{figure}

\subsection{Laplacianos de Hodge y homolog\'ia}
\label{sec:gauge:hodge}

\begin{definition}[Laplaciano de Hodge]
\label{def:gauge:hodge-laplacian}
El \emph{laplaciano de Hodge} de orden $k$ es
\[
  L_k = \partial_{k+1}\partial_{k+1}^{\!\top} + \partial_k^{\!\top}\partial_k
      = L_k^{\text{up}} + L_k^{\text{down}},
\]
con $L_k^{\text{down}} = \partial_k^{\!\top}\partial_k$ el flujo desde las celdas
$(k-1)$ y $L_k^{\text{up}} = \partial_{k+1}\partial_{k+1}^{\!\top}$ el flujo hacia
las celdas $(k+1)$. Para $k = 0$ se reduce al laplaciano del grafo.
\end{definition}

\begin{theorem}[Propiedades del laplaciano de Hodge]
\label{thm:gauge:hodge-props}
$L_k$ es sim\'etrico y semidefinido positivo; su n\'ucleo es isomorfo a la
homolog\'ia $k$-\'esima, $\ker(L_k) \cong H_k(\mathcal{K})$; y $C_k$ se descompone
ortogonalmente como
\[
  C_k = \Img(\partial_{k+1}) \oplus \ker(L_k) \oplus \Img(\partial_k^{\!\top}).
\]
\end{theorem}

\begin{proof}
La simetr\'ia y la positividad son inmediatas, ya que $L_k$ es una suma de
t\'erminos $A^{\!\top} A$ y $\langle L_k\mathbf{x}, \mathbf{x}\rangle =
\|\partial_k\mathbf{x}\|^2 + \|\partial_{k+1}^{\!\top}\mathbf{x}\|^2 \ge 0$. Por
tanto $L_k\mathbf{x} = 0$ si y solo si $\partial_k\mathbf{x} = 0$ y
$\partial_{k+1}^{\!\top}\mathbf{x} = 0$, es decir, $\mathbf{x}$ es un $k$-ciclo
ortogonal a $\Img(\partial_{k+1})$; este es exactamente el espacio de
representantes arm\'onicos de $H_k = \ker\partial_k / \Img\partial_{k+1}$. Los
tres subespacios son mutuamente ortogonales porque $\partial_k\partial_{k+1} = 0$
(\Cref{thm:gauge:partial-squared}), y $\ker L_k$ es el complemento ortogonal de
$\Img(\partial_{k+1}) \oplus \Img(\partial_k^{\!\top})$ en $C_k$.
\end{proof}

La dimensi\'on $\beta_k = \dim H_k(\mathcal{K}) = \operatorname{rank}(Z_k) -
\operatorname{rank}(B_k)$ es el $k$-\'esimo \emph{n\'umero de Betti}, que cuenta
los huecos de dimensi\'on $k$: $\beta_0$ componentes conexas, $\beta_1$ lazos
independientes, $\beta_2$ cavidades. Para la esfera $(\beta_0, \beta_1, \beta_2) =
(1, 0, 1)$ y para el toro $(1, 2, 1)$, dando caracter\'isticas de Euler
$\chi(S^2) = 2$ y $\chi(T^2) = 0$ a trav\'es de la suma alternada $\chi = \sum_k
(-1)^k \beta_k$. En variedades suaves los mismos invariantes surgen
anal\'iticamente: la \emph{cohomolog\'ia de de Rham}
$H^k_\dR(\Mani) = \ker(\mathrm{d}\colon \Omega^k \to \Omega^{k+1}) /
\Img(\mathrm{d}\colon \Omega^{k-1} \to \Omega^k)$ es isomorfa a la cohomolog\'ia
singular con coeficientes reales (teorema de de Rham), ligando el c\'alculo
exterior de \Cref{def:gauge:curvature} con la topolog\'ia. El \'algebra exterior
$\bigwedge V = \bigoplus_k \bigwedge^k V$, con producto antisim\'etrico $v \wedge
w = -w \wedge v$, es precisamente $\Cl(V, 0)$, el \'algebra de Clifford de la
forma nula; este es el puente entre las extensiones algebraica y topol\'ogica.

\subsection{Paso de mensajes y atenci\'on simpliciales}
\label{sec:gauge:simplicial-attention}

La estructura de frontera dicta c\'omo intercambian mensajes las celdas. Para un
$k$-s\'implice $\sigma$ hay tres vecindades naturales: los vecinos
\emph{inferiores} (sus caras $(k-1)$, v\'ia $\partial_k$), los vecinos
\emph{superiores} (las celdas $(k+1)$ de las que es cara, v\'ia $\delta_k$), y los
vecinos \emph{horizontales} (otras celdas $k$ que comparten una cara).
\Cref{fig:gauge:simplicial-messages} muestra los tres tipos de mensaje que
alcanzan una arista objetivo.

\begin{figure}[ht]
\centering
\begin{tikzpicture}[>=Stealth, font=\small, scale=0.78, transform shape,
  msg/.style={line width=1.6pt, -{Stealth[length=6pt, width=4pt]}}]
\coordinate (A) at (0,0); \coordinate (B) at (5,0); \coordinate (C) at (2.5,3.5);
\coordinate (D) at (-3.5,1.2); \coordinate (E) at (-3.0,-2.0); \coordinate (F) at (8.5,1.2);
\fill[gdlorange!15] (A)--(B)--(C)--cycle;
\draw[gdlorange!40] (A)--(B)--(C)--cycle;
\node[gdlorange!70!black] at (2.5,1.6) {$\tau^2$};
\draw[gdlgray!50, line width=1.2pt] (A)--(C) (B)--(C);
\draw[gdlpurple!30, line width=2.2pt] (A)--(D) (A)--(E) (B)--(F);
\draw[gdlred!70, line width=4pt] (A)--(B);
\node[gdlred!80!black, font=\bfseries] at (2.5,-0.55) {$\sigma^1$ (objetivo)};
\foreach \v/\lab/\pos in {A/$v_0$/{below left}, B/$v_1$/{below right}, C/$v_2$/above,
  D/$v_3$/left, E/$v_4$/{below left}, F/$v_5$/right} {
  \fill[gdlblue!60] (\v) circle (5pt);
  \node[font=\footnotesize, gdlblue!70!black, \pos=4pt] at (\v) {\lab};}
\draw[msg, gdlgreen!70] ($(A)+(-0.3,-1.8)$) to[bend right=25] ($(A)+(0.5,-0.25)$);
\draw[msg, gdlgreen!70] ($(B)+(0.3,-1.8)$) to[bend left=25] ($(B)+(-0.5,-0.25)$);
\node[gdlgreen!70!black, font=\scriptsize] at (2.5,-1.7) {$m^{\downarrow}$ (frontera)};
\draw[msg, gdlorange!70] (2.5,2.8) -- (2.5,0.3);
\node[gdlorange!70!black, font=\scriptsize, anchor=south west] at (3.0,2.3) {$m^{\uparrow}$ (cofrontera)};
\draw[msg, gdlpurple!70] ($(A)!0.5!(D)$) to[bend right=20] ($(A)+(0.6,0.35)$);
\draw[msg, gdlpurple!70] ($(A)!0.5!(E)$) to[bend left=20] ($(A)+(0.6,-0.35)$);
\draw[msg, gdlpurple!70] ($(B)!0.5!(F)$) to[bend left=20] ($(B)+(-0.6,0.35)$);
\node[gdlpurple!70!black, font=\scriptsize, anchor=east] at (-3.8,0) {$m^{\leftrightarrow}$ (adyacente)};
\end{tikzpicture}
\caption[Mensajes simpliciales]{Tres tipos de mensaje que alcanzan una arista
objetivo $\sigma^1$ (rojo): mensajes \emph{inferiores} desde sus v\'ertices de
frontera (verde, v\'ia $\partial$), mensajes \emph{superiores} desde los
tri\'angulos de los que es cara (naranja, v\'ia $\delta$), y mensajes
\emph{horizontales} desde aristas adyacentes (p\'urpura).}
\label{fig:gauge:simplicial-messages}
\end{figure}

\begin{definition}[Atenci\'on simplicial]
\label{def:gauge:simplicial-attention}
Una red de atenci\'on simplicial~\citep{giusti2022simplicial} actualiza la
representaci\'on $\mathbf{h}_\sigma$ de un $k$-s\'implice agregando los tres tipos
de mensaje,
\[
  \mathbf{h}_\sigma^{(\ell+1)} = \sigma\!\Big(\mathbf{W}^{\text{self}}
  \mathbf{h}_\sigma^{(\ell)} + \mathbf{m}_\sigma^{\downarrow} +
  \mathbf{m}_\sigma^{\uparrow} + \mathbf{m}_\sigma^{\leftrightarrow}\Big),
  \qquad
  \mathbf{m}_\sigma^{\bullet} = \sum_\tau \alpha_{\tau\to\sigma}\,
  \mathbf{W}^{\bullet}\mathbf{h}_\tau,
\]
con coeficientes de atenci\'on
\[
  \alpha_{\tau\to\sigma} = \softmax_\tau\Big(\mathrm{LeakyReLU}\big(\mathbf{a}^{\!\top}
  [\mathbf{W}_\sigma\mathbf{h}_\sigma \,\|\, \mathbf{W}_\tau\mathbf{h}_\tau \,\|\,
  \mathbf{e}_{\text{orient}}(\tau, \sigma)]\big)\Big)
\]
que incorporan la orientaci\'on relativa $\mathbf{e}_{\text{orient}}$. Plegar un
signo $s_{\tau\to\sigma} \in \{\pm 1\}$ dentro de los coeficientes produce
atenci\'on \emph{orientada} (con signo), preservando la estructura
cohomol\'ogica.
\end{definition}

\begin{theorem}[Equivarianza simplicial]
\label{thm:gauge:simplicial-equivariance}
Las redes de atenci\'on simplicial son equivariantes bajo automorfismos
simpliciales: si $\phi\colon \mathcal{K} \to \mathcal{K}$ es un automorfismo,
entonces $\mathrm{SAN}(\phi(\mathbf{X})) = \phi(\mathrm{SAN}(\mathbf{X}))$.
\end{theorem}

\begin{proof}
Un automorfismo $\phi$ induce una permutaci\'on $P_k$ de los $k$-s\'implices que
preserva la incidencia, $B_k P_{k-1} = P_k B_k$, conmutando por tanto con los
laplacianos, $P_k L_k P_k^{\!\top} = L_k$. La actualizaci\'on depende de las
celdas solo a trav\'es de estas matrices de incidencia y a trav\'es de
coeficientes de atenci\'on que son funciones sim\'etricas de las
representaciones, de modo que $P_k\cdot\mathrm{SAN}(\mathbf{X}) =
\mathrm{SAN}(P_k\cdot\mathbf{X})$ en cada dimensi\'on $k$.
\end{proof}

El laplaciano de Hodge tambi\'en sirve como filtro espectral para la atenci\'on.
Con la descomposici\'on en autovalores $L_k = U_k\Lambda_k U_k^{\!\top}$
(multiplicidad del autovalor cero igual a $\beta_k$, con autovectores
arm\'onicos), se construyen consultas, claves y valores a partir de filtros
espectrales $\mathbf{g}_\theta(L_k) = \sum_i \theta_i T_i(\tilde L_k)$ en la base
de Chebyshev, dando una \emph{atenci\'on por laplaciano de Hodge} que separa las
componentes irrotacional, solenoidal y arm\'onica de las se\~nales de arista. La
\emph{homolog\'ia persistente} a\~nade una dimensi\'on multiescala: sobre una
filtraci\'on $\mathcal{K}_0 \subseteq \cdots \subseteq \mathcal{K}_m$ el diagrama
de persistencia registra el nacimiento $b$ y la muerte $d$ de cada
caracter\'istica homol\'ogica, y los valores de persistencia
$\mathrm{pers}(i, j)$ pueden sesgar los logits de atenci\'on,
$\alpha_{ij} = \softmax(\mathbf{q}_i^{\!\top}\mathbf{k}_j / \sqrt{d} + w\,
\mathrm{pers}(i,j))$, aportando robustez al ruido y una conexi\'on con la
topolog\'ia de los datos. Por \'ultimo, los \emph{haces celulares} (cellular
sheaves) adjuntan un espacio vectorial (tallo) a cada celda y una aplicaci\'on de
restricci\'on a cada incidencia, y el laplaciano del haz
$L_\mathcal{F} = \delta^{\!\top}\delta$ generaliza el laplaciano del grafo a
espacios de caracter\'isticas heterog\'eneos, sustentando las redes neuronales de
haces~\citep{hansen2021sheaf, bodnar2023cellular}. A lo largo de todas estas
construcciones recurre el principio del paso de mensajes, ahora respetando la
incidencia y la orientaci\'on en lugar de la mera adyacencia.

\section{Codificaci\'on posicional, rompimiento de simetr\'ia y extensiones
cu\'anticas}
\label{sec:gauge:positional}

Dos ingredientes adicionales refinan las arquitecturas geom\'etricas: c\'omo
codificar la posici\'on en un dominio, y c\'omo relajar la equivarianza cuando los
datos lo exigen.

\paragraph{Caracter\'isticas de Fourier y el teorema de Bochner.}
Un kernel definido positivo invariante por traslaci\'on admite una
representaci\'on espectral, el fundamento de las caracter\'isticas aleatorias de
Fourier.

\begin{theorem}[Bochner]
\label{thm:gauge:bochner}
Un kernel continuo $k(\mathbf{x}, \mathbf{y}) = k(\mathbf{x} - \mathbf{y})$ en
$\mathbb{R}^d$ es definido positivo si y solo si es la transformada de Fourier de
una medida finita no negativa $\mu$:
\[
  k(\mathbf{x} - \mathbf{y}) = \int_{\mathbb{R}^d} e^{2\pi i\,
  \omega^{\!\top}(\mathbf{x} - \mathbf{y})}\, \mathrm{d}\mu(\omega).
\]
\end{theorem}

Muestrear frecuencias $\omega_j$ de la densidad asociada produce el mapa de
caracter\'isticas $z(\mathbf{x}) = \sqrt{2/m}\,[\cos(2\pi\omega_j^{\!\top}\mathbf{x}),
\sin(2\pi\omega_j^{\!\top}\mathbf{x})]_j$ con $\mathbb{E}[z(\mathbf{x})^{\!\top}
z(\mathbf{y})] = k(\mathbf{x}, \mathbf{y})$ y error uniforme $O(1/\sqrt{m})$. En
una variedad riemanniana compacta las sinusoides se reemplazan por
autofunciones de Laplace--Beltrami $\{\phi_k\}$: el mapa $\gamma_\Mani(p) =
\sqrt{2/m}\,(\phi_1(p), \ldots, \phi_m(p))$ aproxima el kernel de difusi\'on
$k_t(p, q) = \sum_k e^{-t\lambda_k}\phi_k(p)\phi_k(q)$ con error
$O(e^{-t\lambda_{m+1}})$, el an\'alogo intr\'inseco de las caracter\'isticas de
Fourier. En un fibrado principal estas pueden hacerse gauge-equivariantes
transport\'andolas, $\gamma_{\text{gauge}}(p) =
\Gamma_\gamma^{-1}\gamma(\pi(p))$, de modo que la codificaci\'on se transforme
correctamente bajo cambios locales de marco.

\paragraph{Rompimiento controlado de simetr\'ia.}
La equivarianza exacta es a veces demasiado r\'igida: las simetr\'ias reales
pueden ser aproximadas, o las direcciones privilegiadas pueden portar
informaci\'on. \citet{vanderouderaa2024symmetry} interpolan entre una aplicaci\'on
equivariante $\mathcal{F}_{\text{eq}}$ y una sin restricciones,
\[
  \mathcal{F}_\lambda = (1 - \lambda)\,\mathcal{F}_{\text{eq}} + \lambda\,
  \mathcal{F}_{\text{in}},
  \qquad \lambda \in [0, 1].
\]
En $\lambda = 0$ la equivarianza es exacta; conforme $\lambda$ crece se introduce
un sesgo direccional controlado, con $\lambda$ aprendido y regularizado hacia cero
de modo que la red recae por defecto en la equivarianza salvo que los datos
justifiquen romperla.

\begin{proposition}[Transici\'on suave]
\label{prop:gauge:symmetry-breaking}
Si $f_0$ es $G$-equivariante y $h$ arbitraria, entonces $f_\lambda = (1 - \lambda)
f_0 + \lambda h$ satisface $\|f_\lambda(g\cdot x) - \rho(g) f_\lambda(x)\| =
\lambda\,\|h(g\cdot x) - \rho(g) h(x)\|$, que tiende a $0$ uniformemente en
compactos cuando $\lambda \to 0$.
\end{proposition}

\begin{proof}
Desarrollando y usando la equivarianza exacta de $f_0$, los t\'erminos en $f_0$ se
cancelan, dejando $\lambda\,(h(g\cdot x) - \rho(g) h(x))$. En un compacto
$K \subset \Mani \times G$ la funci\'on continua $(x, g) \mapsto \|h(g\cdot x) -
\rho(g) h(x)\|$ est\'a acotada por alg\'un $M$, de modo que la desviaci\'on es a
lo sumo $\lambda M \to 0$.
\end{proof}

Estos principios se extienden incluso a los circuitos cu\'anticos, donde una red
cu\'antica gauge-equivariante es un unitario parametrizado $U(\theta) =
\exp(i\sum_j \theta_j H_j^{\text{gauge}})$ construido a partir de hamiltonianos
que conmutan con los generadores gauge, reflejando la construcci\'on cl\'asica al
nivel de la evoluci\'on unitaria.

\section{S\'intesis: los gauges en la imagen unificada}
\label{sec:gauge:synthesis}

El marco gauge completa el arco de las cinco G, y el mismo hilo de paso de
mensajes que recorri\'o las cuadr\'iculas, los grupos, los grafos y las
geod\'esicas lo recorre tambi\'en. Una convoluci\'on geom\'etrica es una capa
estructurada de paso de mensajes: el grafo subyacente refleja la geometr\'ia del
dominio, las funciones de mensaje preservan la equivarianza al grupo de
simetr\'ia, la agregaci\'on respeta la geometr\'ia local (la convoluci\'on se
convierte en una suma ponderada localmente), y los pesos se comparten a lo largo
de las \'orbitas del grupo. La convoluci\'on gauge-equivariante es este principio
sobre un fibrado principal, con el transporte paralelo suministrando las funciones
de mensaje que alinean los marcos vecinos.

Vistas as\'i, las arquitecturas de todo el libro son facetas de una sola
construcci\'on. Las convoluciones cl\'asicas y de grupo son el caso plano y
globalmente trivializado; el paso de mensajes en grafos y geod\'esicas relaja el
dominio a uno discreto o curvo; la convoluci\'on gauge a\~nade el marco local y su
conexi\'on; y la atenci\'on --- incluyendo las variantes de Clifford y simplicial
de este cap\'itulo --- es paso de mensajes con pesos adaptativos dependientes del
contenido. La transici\'on es incluso cuantitativa: una convoluci\'on
geom\'etrica surge como el l\'imite de temperatura aguda de un kernel de
atenci\'on,
\[
  (f * g)(x) = \lim_{\beta \to \infty} \int_\Mani
  \frac{e^{-\beta d(x, y)^2}}{\int_\Mani e^{-\beta d(x, y')^2}\, \mathrm{d}y'}\,
  g(y)\, f(y)\, \mathrm{d}y,
\]
recuperando una convoluci\'on localizada a partir de la atenci\'on suave cuando
$\beta \to \infty$. \Cref{ch:applications} retoma los beneficios concretos ---
mallas, esferas y mol\'eculas $3$D --- y \Cref{ch:conclusion} re\'une la visi\'on
unificada de paso de mensajes de todo el libro, de la que la imagen gauge es la
expresi\'on m\'as general.
